
% ==================== front ====================
% Самодостаточный исходник. Компиляция: pdflatex near-circular-transfer.tex
% Для оглавления и перекрёстных ссылок требуется повторный запуск pdflatex.
\documentclass[a4paper,11pt]{article}
\usepackage{iftex}
\RequirePDFTeX
\usepackage{cmap}
\usepackage[T2A]{fontenc}
\usepackage[utf8]{inputenc}
\usepackage[english,russian]{babel}
\usepackage{PTSerif}
\usepackage{amsmath,amssymb,amsthm,mathtools}
\usepackage[a4paper,left=20mm,right=20mm,top=22mm,bottom=23mm,headheight=14pt]{geometry}
\usepackage{microtype}
\usepackage{graphicx,xcolor,tikz,booktabs,tabularx,array,longtable}
\usetikzlibrary{arrows.meta,calc}
\usepackage{enumitem,fancyhdr,needspace}
\usepackage{hyperref}
\usepackage{bookmark}
\definecolor{accent}{HTML}{187AA5}
\definecolor{muted}{HTML}{5A6166}
\definecolor{figink}{HTML}{25323F}
\definecolor{figgray}{HTML}{687582}
\definecolor{figrad}{HTML}{176AA6}
\definecolor{figtan}{HTML}{087F78}
\definecolor{fignormal}{HTML}{B57714}
\definecolor{figplane}{HTML}{EEF4F8}
\definecolor{figorbit}{HTML}{79539C}
\hypersetup{unicode=true,colorlinks=true,linkcolor=accent,citecolor=accent,urlcolor=accent,
 pdftitle={Задача перехода в околокруговом приближении},
 pdfauthor={Материалы курса «Введение в маневрирование КА»},
 pdfsubject={Уравнения движения, условия оптимальности, алгоритмы импульсного перехода}}
\newcommand{\dd}{\mathop{}\!\mathrm d}
\newcommand{\R}{\mathbb R}
\DeclareMathOperator{\sgn}{sgn}
\numberwithin{equation}{section}
\newtheorem{theorem}{Теорема}[section]
\newtheorem{lemma}[theorem]{Лемма}
\newtheorem{proposition}[theorem]{Утверждение}
\theoremstyle{remark}
\newtheorem{remark}[theorem]{Замечание}
\renewcommand{\proofname}{Доказательство}
\setcounter{tocdepth}{2}
\setcounter{secnumdepth}{3}
\setlength{\parindent}{1.3em}
\setlength{\parskip}{2pt}
\linespread{1.055}
\setlist{itemsep=2pt,topsep=4pt,parsep=0pt}
\setlength{\emergencystretch}{1.5em}
\clubpenalty=10000
\widowpenalty=10000
\displaywidowpenalty=10000
\allowdisplaybreaks[1]
\pagestyle{fancy}
\fancyhf{}
\fancyhead[L]{\small\color{muted}Задача перехода в околокруговом приближении}
\fancyfoot[C]{\small\thepage}
\renewcommand{\headrulewidth}{.3pt}
\renewcommand{\footrulewidth}{0pt}
\begin{document}
\begin{titlepage}
\centering
{\large Московский физико-технический институт\par}
\vspace{8mm}
{\large Курс «Введение в маневрирование КА»\par}
\vspace{32mm}
{\color{accent}\rule{28mm}{.8pt}\par}
\vspace{9mm}
{\LARGE\bfseries Задача перехода\\[3mm]
в околокруговом приближении\par}
\vspace{10mm}
{\large Методические материалы\par}
\vspace{8mm}
{Уравнения движения. Условия оптимальности.\\
Компланарные и некомпланарные импульсные переходы.\par}
\vfill
{Первая редакция\par}
\vspace{3mm}
{26 сентября 2026 года\par}
\end {titlepage}
\setcounter{page}{2}
\tableofcontents
\clearpage

\section*{Постановка задачи и обозначения}
\addcontentsline{toc}{section}{Постановка задачи и обозначения}

Рассматривается переход космического аппарата между двумя близкими
орбитами малого эксцентриситета в центральном ньютоновском поле.
Начальная и конечная орбиты заданы; положение аппарата на конечной
орбите не назначается. Начало маневрирования можно выбрать в любой
точке начальной орбиты. Время перехода свободно.
Управление представляет собой последовательность мгновенных
изменений скорости. Требуется определить число импульсов,
места их приложения, направления и величины так, чтобы минимизировать
\[
 \Delta V_\Sigma=\sum_{j=1}^N\lVert\Delta\boldsymbol V_j\rVert.
\]
Это характеристическая скорость, а не сумма векторов скоростей
и не механический импульс $m\Delta\boldsymbol V$.
При постоянной эффективной скорости истечения $c$ и отсутствии
сброса массы между манёврами уравнение Циолковского даёт
$m_f/m_i=\exp(-\Delta V_\Sigma/c)$.
Поэтому минимизация характеристической скорости равносильна
минимизации расхода топлива при фиксированной начальной массе.

Все орбиты перехода рассматриваются в малой окрестности одного
опорного круга. Малыми считаются эксцентриситеты, относительные
отклонения больших полуосей, взаимные наклонения и отношения
импульсов скорости к круговой скорости. В уравнениях сохраняется
первый порядок по этим величинам. Ограничения на длительность
перехода, высоту промежуточной орбиты, число включений двигателя,
направления тяги и доступные участки орбиты здесь не задаются.
При введении таких ограничений задача оптимизации изменяется.

Изложение опирается на теорию базис-вектора Лоудена
\cite{lawden,lawden-ru}, решение Эдельбаума для близких околокруговых
орбит \cite{edelbaum}, линейную часть отчёта \cite{gobetz}
и вторую главу работ Баранова \cite{baranov-ru,baranov-en}.
Орбитальные обозначения и нормировки приняты преимущественно
по Баранову. Крышка над координатами и скоростями обозначает их нормировку:
длины делятся на $r_0$, скорости --- на $V_0$.
Для импульсов $\Delta\widehat{\boldsymbol V}_j=\Delta\boldsymbol V_j/V_0$.
Все промежуточные преобразования выполняются в единой системе осей.

\begin{center}
\renewcommand{\arraystretch}{1.08}
\begin{tabularx}{\textwidth}{@{}>{$}l<{$}X@{}}
\toprule
\multicolumn{1}{@{}l}{Обозначение}&Смысл\\
\midrule
 \mu & Гравитационный параметр центрального тела.\\
 a_i,a_f & Большие полуоси начальной и конечной орбит.\\
 r_0=a_0 & Радиус опорной круговой орбиты; обычно $(a_i+a_f)/2$.\\
 V_0,n_0 & Круговая скорость и среднее движение:
 $V_0=\sqrt{\mu/r_0}$, $n_0=V_0/r_0$.\\
 r,u,z & Цилиндрические координаты относительно начальной плоскости.\\
 \varphi,\theta & Абсолютная фаза опорного круга и угол, прошедший
 после отдельного импульса.\\
 \Delta a & Безразмерное изменение большой полуоси $(a_f-a_i)/r_0$.\\
 \Delta e_x,\ \Delta e_y & Компоненты разности векторов эксцентриситета
 в общих неподвижных осях.\\
 \Delta i & Малый взаимный угол между плоскостями; ${\Delta i}\geqslant0$.\\
 \Delta\widehat{\boldsymbol V}_j & Безразмерный вектор $j$-го импульса, то есть
 $\Delta\boldsymbol V_j/V_0$ в локальных осях.\\
 \Delta\widehat V_j,\ J & Модуль отдельного импульса и их сумма:
 ${\Delta\widehat V_j}=\lVert\Delta\widehat{\boldsymbol V}_j\rVert$, $J=\sum_j{\Delta\widehat V_j}$.\\
 \boldsymbol l,\boldsymbol p & Множители граничных условий и базис-вектор.\\
\bottomrule
\end{tabularx}
\end{center}

Углы в формулах измеряются в радианах; в расчётных примерах
места импульсов дополнительно приводятся в градусах.
\clearpage

% ==================== dynamics ====================
\section{Уравнения движения}
\label{sec:dynamics}

\subsection{Исходная модель и выбор координат}

Рассматривается движение материальной точки в центральном ньютоновском
поле. Гравитационный параметр притягивающего тела обозначается через
$\mu$. Нецентральность поля, сопротивление атмосферы и внешние
гравитационные возмущения в рассматриваемую модель не входят.
Управляющее ускорение может быть произвольным по направлению; ниже
из уравнений с конечным ускорением будет получена импульсная модель.

Пусть $\boldsymbol R$ --- радиус-вектор аппарата относительно центра
притяжения, $\boldsymbol V$ --- его скорость в инерциальной системе
координат, $\boldsymbol a_{\text{упр}}$ --- управляющее ускорение.
Исходные уравнения имеют вид
\begin {equation}
 \dot{\boldsymbol R}=\boldsymbol V,\qquad
 \dot{\boldsymbol V}=-\frac{{\mu}}{{\lVert\boldsymbol R\rVert^3}}
 \boldsymbol R+\boldsymbol a_{\text{упр}}.
 \label {eq:cartesian-dynamics}
\end {equation}
Точкой обозначается производная по физическому времени $t$.

За неподвижную опорную плоскость $OXY$ примем плоскость начальной орбиты.
Начало координат $O$ находится в центре притяжения. Ось $Z$ направим
по нормали к этой плоскости так, чтобы положительное вращение вокруг
неё соответствовало движению по начальной орбите.
Введём цилиндрические координаты $(r,u,z)$: $r$ --- расстояние до
нормальной оси, $u$ --- угол в опорной плоскости, возрастающий в
направлении движения по начальной орбите, $z$ --- высота над этой
плоскостью. При этом расстояние до центра притяжения равно
$s=\sqrt{r^2+z^2}$, а не $r$. Геометрия координат и локального базиса
показана на рис.~\ref{fig:cylindrical-coordinates}. Единичные векторы базиса
записываются в неподвижных декартовых осях как
\begin {equation}
 \boldsymbol e_r=(\cos u,\sin u,0),\qquad
 \boldsymbol e_t=(-\sin u,\cos u,0),\qquad
 \boldsymbol e_z=(0,0,1).
 \label {eq:cylindrical-basis}
\end {equation}

\begin{figure}[!htbp]
\centering
\resizebox{\textwidth}{!}{%
% Единая 3D-проекция: X=(1,-.16), Y=(-.5,.58), Z=(0,1).
% Две круговые орбиты радиуса a=4.3. Наклон текущей орбиты 30 градусов.
% P имеет аргумент широты 30 градусов. V — точная касательная в P.
\begin{tikzpicture}[x=1cm,y=1cm,font=\small,
 line cap=round,line join=round,
 fixed/.style={draw=figgray,line width=.65pt,-{Stealth[length=2.1mm]}},
 local/.style={line width=1.35pt,-{Stealth[length=2.7mm,width=1.8mm]}},
 guide/.style={draw=figgray!60,line width=.6pt,dash pattern=on 2.4pt off 2.4pt},
 point/.style={circle,inner sep=2pt,fill=figink},
 lab/.style={fill=white,inner sep=1.5pt}]
\path[use as bounding box] (0,-.05) rectangle (15.9,11.2);
\pgfmathsetmacro{\orba}{4.3}
\pgfmathsetmacro{\orbi}{30}
\pgfmathsetmacro{\orbv}{30}
\pgfmathsetmacro{\px}{\orba*cos(\orbv)}
\pgfmathsetmacro{\py}{\orba*sin(\orbv)*cos(\orbi)}
\pgfmathsetmacro{\pz}{\orba*sin(\orbv)*sin(\orbi)}
\pgfmathsetmacro{\phiu}{atan2(\py,\px)}
\begin{scope}[shift={(7.5,5.2)},
 x={(1cm,-.16cm)},y={(-.5cm,.58cm)},z={(0cm,1cm)}]
 \coordinate (O) at (0,0,0);
 \coordinate (H) at (\px,\py,0);
 \coordinate (P) at (\px,\py,\pz);
 \filldraw[fill=figplane,draw=figgray!30,line width=.6pt]
   (-4.7,-4.7,0)--(4.7,-4.7,0)--(4.7,4.7,0)--(-4.7,4.7,0)--cycle;
 % Начальная орбита определяет опорную плоскость OXY.
 \draw[draw=figgray!85,line width=.9pt,dash pattern=on 4pt off 2.5pt]
   plot[domain=0:360,samples=180,variable=\t]
   ({\orba*cos(\t)},{\orba*sin(\t)},0);
 \draw[guide] (-1,0,0)--(O);
 \draw[guide] (0,-.8,0)--(O);
 \draw[fixed] (O)--(5.8,0,0) node[below right=2pt] {$X$};
 \draw[fixed] (O)--(0,5.3,0) node[above left=2pt] {$Y$};
 \draw[fixed] (O)--(0,0,1.15) node[above=3pt] {$Z$};
 % Текущая орбита: нижняя полудуга находится под опорной плоскостью.
 \draw[draw=figorbit!75,line width=1pt,dash pattern=on 4pt off 2.5pt]
   plot[domain=180:360,samples=100,variable=\t]
   ({\orba*cos(\t)},{\orba*sin(\t)*cos(\orbi)},
    {\orba*sin(\t)*sin(\orbi)});
 \draw[draw=figorbit,line width=1.3pt]
   plot[domain=0:180,samples=100,variable=\t]
   ({\orba*cos(\t)},{\orba*sin(\t)*cos(\orbi)},
    {\orba*sin(\t)*sin(\orbi)});
 \draw[draw=figorbit,line width=1.3pt,-{Stealth[length=2.5mm]}]
   plot[domain=135:154,samples=20,variable=\t]
   ({\orba*cos(\t)},{\orba*sin(\t)*cos(\orbi)},
    {\orba*sin(\t)*sin(\orbi)});
 \coordinate (orbitlabel) at ({\orba*cos(115)},
   {\orba*sin(115)*cos(\orbi)},{\orba*sin(115)*sin(\orbi)});
 \coordinate (initiallabel) at ({\orba*cos(300)},{\orba*sin(300)},0);
 % r и u находятся в опорной плоскости; z — вдоль её нормали.
 \path[fill=figrad!10] (O)--(1.4,0,0)
   plot[domain=0:\phiu,samples=25,variable=\t]
   ({1.4*cos(\t)},{1.4*sin(\t)},0)--cycle;
 \draw[draw=figrad,line width=1pt] (O)--(H)
   node[pos=.63,above=4pt,text=figrad,lab] {$r$};
 \draw[guide] (H)--(P);
 \draw[draw=figink,line width=1.2pt,-{Stealth[length=2.6mm]}]
   (O)--(P) node[pos=.60,above left=3pt,lab] {$\boldsymbol R$};
 \draw[draw=figrad,line width=.9pt,-{Stealth[length=1.8mm]}]
   plot[domain=0:\phiu,samples=25,variable=\t]
   ({1.4*cos(\t)},{1.4*sin(\t)},0);
 \node[text=figrad] at ({1.83*cos(12)},{1.83*sin(12)},0) {$u$};
 % Положительные направления базиса, единичная длина в одном масштабе.
 \draw[local,figrad] (P)--++({1.8*cos(\phiu)},{1.8*sin(\phiu)},0)
   node[above right=4pt,align=left,text=figrad]
   {$\boldsymbol e_r$\\[2pt]$V_r=\dot r,\quad S$};
 \draw[local,figtan] (P)--++({-1.8*sin(\phiu)},{1.8*cos(\phiu)},0)
   node[above left=4pt,xshift=-3mm,text=figtan]
   {$\boldsymbol e_t,\quad V_t=r\dot u,\quad T$};
 \draw[local,fignormal] (P)--++(0,0,1.8)
   node[above=4pt,align=center,text=fignormal]
   {$\boldsymbol e_z$\\[2pt]$V_z=\dot z,\quad W$};
 % Вектор скорости направлен по касательной к текущей орбите.
 \draw[draw=figorbit,line width=1.55pt,-{Stealth[length=2.8mm,width=1.9mm]}]
   (P)--++({-1.8*sin(\orbv)},
   {1.8*cos(\orbv)*cos(\orbi)},{1.8*cos(\orbv)*sin(\orbi)})
   node[above left=3pt,text=figorbit,lab] {$\boldsymbol V$};
 \node[point] at (O) {};
 \node[below left=3pt] at (O) {$O$};
 \node[point,fill=figrad] at (H) {};
 \node[below right=4pt,text=figrad,lab] at (H) {$H$};
 \node[point,fill=figorbit] at (P) {};
 \node[right=4pt,xshift=4mm,yshift=-4mm,text=figorbit] at (P) {$P$};
\end{scope}

\coordinate (HZ) at ($(H)+(3.05,0)$);
\coordinate (PZ) at ($(P)+(3.05,0)$);
\draw[guide] ($(H)+(.14,0)$)--($(HZ)+(.2,0)$);
\draw[guide] ($(P)+(.14,0)$)--($(PZ)+(.2,0)$);
\draw[fignormal,line width=.85pt,{Stealth[length=1.8mm]}-{Stealth[length=1.8mm]}]
 (HZ)--(PZ) node[midway,right=4pt,text=fignormal] {$z$};
\draw[figgray,line width=.7pt]
 ($(H)+(-.23,-.040)$)--($(H)+(-.23,.18)$)--($(H)+(0,.22)$);

\node[anchor=west,text=figorbit,font=\small\bfseries] at (1.7,10.05)
 {Орбита КА};
\draw[draw=figorbit!70,line width=.6pt] (3.15,9.88)--(orbitlabel);
\node[anchor=west,text=figgray,align=left] at (13.15,2.9) {Начальная\\орбита};
\draw[draw=figgray!70,line width=.6pt] (13.05,2.9)--(initiallabel);
\node[text=figgray,font=\footnotesize] at (7.3,4.10) {Плоскость $OXY$};
\node[text=figink] at (7.6,.18)
 {$\boldsymbol R=r\boldsymbol e_r+z\boldsymbol e_z,\qquad
 s=\lVert\boldsymbol R\rVert=\sqrt{r^2+z^2}$};
\end{tikzpicture}
}
\caption{Цилиндрические координаты и локальный базис.
Аппарат $P$ находится на наклонной орбите; $\boldsymbol V$ направлен
по касательной, стрелка на орбите задаёт направление движения.
$H$ --- проекция $P$ на плоскость начальной орбиты $OXY$.
Векторы $\boldsymbol e_r,\boldsymbol e_t,\boldsymbol e_z$ задают
положительные направления компонент скорости $V_r,V_t,V_z$
и управляющего ускорения $S,T,W$ соответственно.
Для иллюстрации обе орбиты изображены круговыми; часть наклонной
орбиты ниже плоскости $OXY$ показана фиолетовым пунктиром.}
\label{fig:cylindrical-coordinates}
\end{figure}
Соответственно,
\[
 \dot{\boldsymbol e}_r=\dot u\boldsymbol e_t,\qquad
 \dot{\boldsymbol e}_t=-\dot u\boldsymbol e_r,
 \qquad \dot{\boldsymbol e}_z=0.
\]
Разложим управление:
\[
 \boldsymbol a_{\text{упр}}
 =S\boldsymbol e_r+T\boldsymbol e_t+W\boldsymbol e_z.
\]
Здесь $S$, $T$ и $W$ имеют размерность ускорения. В точной
пространственной задаче $\boldsymbol e_r$ направлен по проекции
радиус-вектора на опорную плоскость, а $\boldsymbol e_z$ не обязан
совпадать с нормалью к текущей орбите. Эти различия следует учитывать
до выполнения линеаризации.

Из равенства $\boldsymbol R=r\boldsymbol e_r+z\boldsymbol e_z$
последовательным дифференцированием получаем
\begin {align*}
 \boldsymbol V&=\dot r\boldsymbol e_r+r\dot u\boldsymbol e_t
                 +\dot z\boldsymbol e_z,\\
 \ddot{\boldsymbol R}
 &=(\ddot r-r\dot u^2)\boldsymbol e_r
 +(r\ddot u+2\dot r\dot u)\boldsymbol e_t
 +\ddot z\boldsymbol e_z.
\end {align*}
Подстановка в \eqref{eq:cartesian-dynamics} даёт точные уравнения
\begin {equation}
 \begin {aligned}
  \ddot r-r\dot u^2&=-\frac{{\mu r}}{{(r^2+z^2)^{3/2}}}+S,\\
  r\ddot u+2\dot r\dot u&=T,\\
  \ddot z&=-\frac{{\mu z}}{{(r^2+z^2)^{3/2}}}+W.
 \end {aligned}
 \label {eq:cylindrical-second-order}
\end {equation}
Удобнее использовать компоненты физической скорости
$V_r=\dot r$, $V_t=r\dot u$ и $V_z=\dot z$. Тогда
\begin {equation}
 \begin {aligned}
  \dot r&=V_r,&
  \dot u&=\frac{{V_t}}{{r}},&
  \dot z&=V_z,\\
  \dot V_r&=\frac{{V_t^2}}{{r}}
       -\frac{{\mu r}}{{(r^2+z^2)^{3/2}}}+S,&
  \dot V_t&=-\frac{{V_rV_t}}{{r}}+T,&
  \dot V_z&=-\frac{{\mu z}}{{(r^2+z^2)^{3/2}}}+W.
 \end {aligned}
 \label {eq:cylindrical-first-order}
\end {equation}
Уравнение для $V_t$ получается дифференцированием произведения
$r\dot u$, поэтому в нём остаётся один член $-V_rV_t/r$, а не два.

\subsection{Линеаризация в окрестности круговой орбиты}

Пусть радиус опорной круговой орбиты равен
\begin {equation}
 r_0=a_0=\frac{{a_i+a_f}}{{2}},\qquad
 V_0=\sqrt{\frac{{\mu}}{{r_0}}},\qquad
 n_0=\sqrt{\frac{{\mu}}{{r_0^3}}}=\frac{{V_0}}{{r_0}}.
 \label {eq:reference-scales}
\end {equation}
Индексы $i$ и $f$ относятся к начальной и конечной орбитам.
Выбор среднего значения больших полуосей удобен, но не является
условием вывода: допустим любой близкий к обеим орбитам опорный круг.
Его фазу зададим как
\[
 \varphi(t)=n_0t+\varphi_0.
\]
Опорное решение имеет вид
$r=r_0$, $V_r=0$, $V_t=V_0$, $z=V_z=0$, $u=\varphi$.
Введём отклонения
\[
 r=r_0+\Delta r,\quad V_r=\Delta V_r,\quad
 V_t=V_0+\Delta V_t,\quad u=\varphi+\Delta u.
\]

При линеаризации координатной системы предполагается, что
\begin {equation}
 \frac{{|\Delta r|}}{{r_0}},\quad
 \frac{{|\Delta V_r|}}{{V_0}},\quad
 \frac{{|\Delta V_t|}}{{V_0}},\quad
 \frac{{|z|}}{{r_0}},\quad
 \frac{{|V_z|}}{{V_0}},\quad |\Delta u|=O(\varepsilon),
 \qquad \varepsilon\ll1.
 \label {eq:state-smallness}
\end {equation}
Управление также считается возмущением первого порядка:
$S,T,W=O(\varepsilon n_0V_0)$ на рассматриваемом интервале.
Для импульсного управления соответствующим условием служит
$\lVert\Delta\boldsymbol V_j\rVert/V_0=O(\varepsilon)$.
Члены второго и более высоких порядков по $\varepsilon$
отбрасываются. В частности,
\begin {align*}
 \frac{{V_t^2}}{{r}}
 &=\frac{{V_0^2}}{{r_0}}
   +\frac{{2V_0}}{{r_0}}\Delta V_t
   -\frac{{V_0^2}}{{r_0^2}}\Delta r+O(\varepsilon^2 n_0V_0),\\
 \frac{{\mu r}}{{(r^2+z^2)^{3/2}}}
 &=\frac{{\mu}}{{r_0^2}}
   -\frac{{2\mu}}{{r_0^3}}\Delta r+O(\varepsilon^2 n_0V_0),\\
 \frac{{V_rV_t}}{{r}}&=n_0\Delta V_r+O(\varepsilon^2 n_0V_0),\qquad
 \frac{{\mu z}}{{(r^2+z^2)^{3/2}}}
 =n_0^2z+O(\varepsilon^2 n_0V_0).
\end {align*}
Второе разложение не содержит линейного члена по $z$: в радиальное
ускорение высота входит через $z^2$. Постоянные члены первого и
второго разложений сокращаются вследствие равенства
$V_0^2/r_0=\mu/r_0^2$. Разложение угловой скорости даёт
\[
 \frac{{V_t}}{{r}}=n_0+\frac{{\Delta V_t-n_0\Delta r}}{{r_0}}
 +O(\varepsilon^2n_0).
\]
Таким образом, система первого приближения принимает вид
\begin {equation}
 \boxed{\begin {aligned}
  \Delta\dot r&=\Delta V_r,\\
  \Delta\dot V_r&=n_0^2\Delta r+2n_0\Delta V_t+S,\\
  \Delta\dot V_t&=-n_0\Delta V_r+T,\\
  \Delta\dot u&=\frac{{\Delta V_t-n_0\Delta r}}{{r_0}},\\
  \dot z&=V_z,\\
  \dot V_z&=-n_0^2z+W.
 \end {aligned}}
 \label {eq:linear-dynamics}
\end {equation}
Здесь и далее равенства для линеаризованной модели записываются без
остаточных членов. Их отсутствие в записи не означает точности
этих равенств для исходной нелинейной задачи.

Связь с обычной формой уравнений относительного движения получается
при введении $\xi=\Delta r$ и $\eta=r_0\Delta u$:
\begin {equation}
 \ddot\xi-2n_0\dot\eta-3n_0^2\xi=S,\qquad
 \ddot\eta+2n_0\dot\xi=T,\qquad
 \ddot z+n_0^2z=W.
 \label {eq:hill-equations}
\end {equation}
Действительно, $\dot\eta=\Delta V_t-n_0\Delta r$.
Поэтому $\dot\eta$ и $\Delta V_t$ --- разные величины: первая
есть производная относительной продольной координаты во вращающейся
системе, вторая --- отклонение физической трансверсальной скорости.

\subsection{Свободное движение и отклик на импульс}

Крышкой будем обозначать нормировку координат на $r_0$, а скоростей
на $V_0$. Безразмерные отклонения и компоненты состояния имеют вид
\begin {equation}
 \Delta\widehat r=\frac{{\Delta r}}{{r_0}},\quad
 \Delta\widehat V_r=\frac{{\Delta V_r}}{{V_0}},\quad
 \Delta\widehat V_t=\frac{{\Delta V_t}}{{V_0}},\quad
 \widehat z=\frac{{z}}{{r_0}},\quad
 \widehat V_z=\frac{{V_z}}{{V_0}}.
 \label {eq:dimensionless-state}
\end {equation}
На пассивном участке $S=T=W=0$. Если штрих обозначает
дифференцирование по $\varphi$, то \eqref{eq:linear-dynamics} даёт
\begin{equation}
 \begin{aligned}
  (\Delta\widehat r)'&=\Delta\widehat V_r,&
  (\Delta u)'&=\Delta\widehat V_t-\Delta\widehat r,\\
  (\Delta\widehat V_r)'&=\Delta\widehat r+2\Delta\widehat V_t,&
  (\widehat z)'&=\widehat V_z,\\
  (\Delta\widehat V_t)'&=-\Delta\widehat V_r,&
  (\widehat V_z)'&=-\widehat z.
 \end{aligned}
 \label{eq:dimensionless-free}
\end{equation}
Из уравнений для $\Delta\widehat r$ и $\Delta\widehat V_t$
следует, что
$\alpha_a=2(\Delta\widehat r+\Delta\widehat V_t)$ постоянно. После исключения $\Delta\widehat V_t$ получаем
$(\Delta\widehat r)''+\Delta\widehat r=\alpha_a$. Поэтому общее свободное движение можно
представить в виде
\begin{equation}
 \begin{aligned}
  h(\varphi)&=e_x\cos\varphi+e_y\sin\varphi,\\
  \Delta\widehat r&=\alpha_a-h(\varphi),\\
  \Delta\widehat V_r&=e_x\sin\varphi-e_y\cos\varphi,\\
  \Delta\widehat V_t&=h(\varphi)-\frac{\alpha_a}{2},\\
  \widehat z&=i_x\sin\varphi-i_y\cos\varphi,\\
  \widehat V_z&=i_x\cos\varphi+i_y\sin\varphi.
 \end{aligned}
 \label{eq:free-elements-solution}
\end{equation}
Здесь $h$ --- вспомогательная скалярная функция, а
$\alpha_a,e_x,e_y,i_x,i_y$ --- пять постоянных интегрирования.
Шестая постоянная задаёт начальную фазу. Если $\varphi_*$ соответствует
моменту $t_*$, то интегрирование $(\Delta u)'=\Delta\widehat V_t-\Delta\widehat r$ даёт
\begin {equation}
 \begin {aligned}
 \Delta u(\varphi)={}&\Delta u(\varphi_*)-
 \frac{3}{2}\alpha_a(\varphi-\varphi_*)\\
 &+2e_x(\sin\varphi-\sin\varphi_*)
 -2e_y(\cos\varphi-\cos\varphi_*).
 \end {aligned}
 \label {eq:free-phase}
\end {equation}
Физический смысл этих постоянных будет установлен в следующем
подразделе.

Идеальный импульс в момент $t_j$ оставляет положение непрерывным и
вызывает скачок скорости:
\begin {equation}
 \boldsymbol R(t_j+0)=\boldsymbol R(t_j-0),\qquad
 \boldsymbol V(t_j+0)=\boldsymbol V(t_j-0)+\Delta\boldsymbol V_j.
 \label {eq:impulse-jump}
\end {equation}
Его локальные безразмерные компоненты обозначим
\begin {equation}
 \Delta\widehat{\boldsymbol V}_j=({\Delta\widehat V_{rj}},{\Delta\widehat V_{tj}},{\Delta\widehat V_{zj}})^{\mathsf T}
 =\frac{{1}}{{V_0}}(\Delta V_{rj},\Delta V_{tj},\Delta V_{zj})^{\mathsf T}.
 \label {eq:normalized-impulse}
\end {equation}
Здесь индекс $j$ относится к отдельному импульсу: $\Delta\widehat V_{rj}$,
$\Delta\widehat V_{tj}$ и $\Delta\widehat V_{zj}$ --- скачки скорости,
делённые на $V_0$. Величины $\Delta\widehat V_r$, $\Delta\widehat V_t$
и $\widehat V_z$ без индекса $j$ описывают текущее состояние аппарата.
При близости к опорному кругу цилиндрический базис отличается от
орбитального базиса текущей орбиты, составленного из радиального,
трансверсального и нормального направлений на углы
порядка $O(\varepsilon)$. Поэтому компоненты малого импульса в этих
базисах различаются лишь на $O(\varepsilon^2V_0)$; в формулах
первого приближения это различие не учитывается.
В силу линейности удобно отдельно рассматривать добавку к движению,
созданную одним импульсом. Непосредственно после него эта добавка
удовлетворяет условиям
\[
 \Delta\widehat r=\Delta u=\widehat z=0,\qquad
 \Delta\widehat V_r={\Delta\widehat V_{rj}},\quad \Delta\widehat V_t={\Delta\widehat V_{tj}},\quad\widehat V_z={\Delta\widehat V_{zj}}.
\]
Введём \emph{прошедший угол}
$\theta=n_0(t-t_j)$. Подстановка начальных условий в
\eqref{eq:dimensionless-free} даёт
\begin {equation}
 \begin {aligned}
  \Delta\widehat r&={\Delta\widehat V_{rj}}\sin\theta+2{\Delta\widehat V_{tj}}(1-\cos\theta),\\
  \Delta\widehat V_r&={\Delta\widehat V_{rj}}\cos\theta+2{\Delta\widehat V_{tj}}\sin\theta,\\
  \Delta\widehat V_t&=-{\Delta\widehat V_{rj}}\sin\theta+{\Delta\widehat V_{tj}}(2\cos\theta-1),\\
  \Delta u&=2{\Delta\widehat V_{rj}}(\cos\theta-1)
                  +{\Delta\widehat V_{tj}}(4\sin\theta-3\theta),\\
  \widehat z&={\Delta\widehat V_{zj}}\sin\theta,\qquad
  \widehat V_z={\Delta\widehat V_{zj}}\cos\theta.
 \end {aligned}
 \label {eq:impulse-response}
\end {equation}
Например, первое выражение получается из
$(\Delta\widehat r)''+\Delta\widehat r=2{\Delta\widehat V_{tj}}$ при $(\Delta\widehat r)(0)=0$ и $(\Delta\widehat r)'(0)={\Delta\widehat V_{rj}}$.
Затем $\Delta\widehat V_r=(\Delta\widehat r)'$, $\Delta\widehat V_t={\Delta\widehat V_{tj}}-\Delta\widehat r$, а $\Delta u$ находится
интегрированием $(\Delta u)'={\Delta\widehat V_{tj}}-2\Delta\widehat r$.

Углы $\theta$ и $\varphi$ решают разные задачи. В
\eqref{eq:impulse-response} $\theta=0$ всегда соответствует моменту
рассматриваемого импульса. Угол $\varphi_j=\varphi(t_j)$ задаёт его
положение относительно единой неподвижной оси. На пассивном участке
$\varphi=\varphi_j+\theta$. Подмена одного угла другим без сдвига начала
отсчёта меняет знаки и ориентацию орбитальных векторов.

В качестве проверки рассмотрим положительный трансверсальный импульс
${\Delta\widehat V_{tj}}>0$, ${\Delta\widehat V_{rj}}={\Delta\widehat V_{zj}}=0$. Через половину опорного периода,
при $\theta=\pi$, получаем
\[
 \Delta r=4r_0{\Delta\widehat V_{tj}},\quad \Delta V_r=0,\quad
 \Delta V_t=-3V_0{\Delta\widehat V_{tj}},\quad \Delta u=-3\pi {\Delta\widehat V_{tj}}.
\]
Аппарат достигает противоположной апсидальной точки возмущённой
орбиты. Отрицательное отклонение трансверсальной скорости в этой
точке согласуется с увеличением радиуса и сохранением момента
количества движения на пассивном участке. Отставание по фазе
связано с увеличением периода. Формула описывает добавку от
импульса; при произвольной начальной орбите к ней прибавляется
свободное движение с исходными начальными условиями.

\subsection{Связь с параметрами орбиты}
\label{sec:orbital-variables}

Для задачи перехода требуется связать движение с геометрией
орбиты. Классические элементы $e,\omega,i,\Omega$ неудобны около
круга и опорной плоскости: при $e=0$ не определён перицентр,
а при $i=0$ --- линия узлов. Векторы эксцентриситета и наклонения
сохраняют определённость и в этих случаях. Ниже их компоненты
выводятся непосредственно из состояния аппарата.

Удельная механическая энергия и большая полуось удовлетворяют
точному равенству
\[
 \mathcal E=\frac{{\lVert\boldsymbol V\rVert^2}}{{2}}
             -\frac{{\mu}}{{s}}=-\frac{{\mu}}{{2a}}.
\]
В первом порядке относительно опорного круга
\[
 \mathcal E+\frac{{V_0^2}}{{2}}=V_0^2(\Delta\widehat r+\Delta\widehat V_t),\qquad
 -\frac{{\mu}}{{2a}}+\frac{{\mu}}{{2r_0}}
       =\frac{{V_0^2}}{{2}}\frac{{a-r_0}}{{r_0}}.
\]
Следовательно,
\begin {equation}
 \alpha_a=\frac{{a-r_0}}{{r_0}}=2(\Delta\widehat r+\Delta\widehat V_t).
 \label {eq:semimajor-state}
\end {equation}
Таким образом, постоянная $\alpha_a$ характеризует отклонение
большой полуоси от опорной. Она отличается от заданного изменения
${\Delta a}=(a_f-a_i)/r_0=\alpha_{a,f}-\alpha_{a,i}$.

Точный вектор эксцентриситета имеет вид
\begin {equation}
 \boldsymbol e=
 \frac{{\boldsymbol V\times(\boldsymbol R\times\boldsymbol V)}}{{\mu}}
 -\frac{{\boldsymbol R}}{{s}}
 =\frac{{(\lVert\boldsymbol V\rVert^2-\mu/s)\boldsymbol R
       -(\boldsymbol R\mathbin{\cdot}\boldsymbol V)\boldsymbol V}}{{\mu}}.
 \label {eq:eccentricity-vector-exact}
\end {equation}
Здесь
$\lVert\boldsymbol V\rVert^2-\mu/s=V_0^2(2\Delta\widehat V_t+\Delta\widehat r)$ и
$\boldsymbol R\mathbin{\cdot}\boldsymbol V=r_0V_0\Delta\widehat V_r$
с точностью первого порядка. Поэтому
\[
 \boldsymbol e=(\Delta\widehat r+2\Delta\widehat V_t)\boldsymbol e_r-\Delta\widehat V_r\boldsymbol e_t
 +O(\varepsilon^2).
\]
Вертикальная компонента отсутствует в первом порядке: она имеет
порядок произведения эксцентриситета на наклонение. При переходе к
неподвижным осям $u$ можно заменить на $\varphi$: ошибка поворота
$O(\Delta u)$ умножается на вектор порядка $O(\varepsilon)$.
Итак,
\begin {equation}
 \boxed{\begin {aligned}
  e_x&=(\Delta\widehat r+2\Delta\widehat V_t)\cos\varphi+\Delta\widehat V_r\sin\varphi,\\
  e_y&=(\Delta\widehat r+2\Delta\widehat V_t)\sin\varphi-\Delta\widehat V_r\cos\varphi.
 \end {aligned}}
 \label {eq:eccentricity-state}
\end {equation}
Для орбиты в опорной плоскости это обычные компоненты
$e_x=e\cos\varpi$, $e_y=e\sin\varpi$, где $\varpi$ отсчитывается
от выбранной неподвижной оси. При $e=0$ обе компоненты равны нулю,
и вводить произвольное значение $\varpi$ не требуется.
Решение \eqref{eq:free-elements-solution} теперь приобретает
непосредственный геометрический смысл:
\[
 \frac{{r-r_0}}{{r_0}}
 =\frac{{a-r_0}}{{r_0}}-e_x\cos\varphi-e_y\sin\varphi.
\]
Это также следует из линеаризации точного полярного уравнения
$r=a(1-e^2)/(1+e\cos f)$.

Для описания плоскости рассмотрим единичную нормаль
$\boldsymbol k=(\boldsymbol R\times\boldsymbol V)/
\lVert\boldsymbol R\times\boldsymbol V\rVert$.
В первом порядке вектор момента равен
\[
 \boldsymbol R\times\boldsymbol V
 =r_0V_0\bigl[(1+\Delta\widehat r+\Delta\widehat V_t)\boldsymbol e_z
             -\widehat z\boldsymbol e_r-\widehat V_z\boldsymbol e_t\bigr].
\]
После нормировки
\[
 \boldsymbol k=\boldsymbol e_z-\widehat z\boldsymbol e_r
                                      -\widehat V_z\boldsymbol e_t.
\]
Введём компоненты наклонения так, чтобы
$\boldsymbol k=(i_y,-i_x,1)+O(\varepsilon^2)$. Получаем
\begin {equation}
 \boxed{\begin {aligned}
  i_x&=\widehat z\sin\varphi+\widehat V_z\cos\varphi,\\
  i_y&=-\widehat z\cos\varphi+\widehat V_z\sin\varphi.
 \end {aligned}}
 \label {eq:inclination-state}
\end {equation}
При малом наклонении $i$ к опорной плоскости и долготе восходящего
узла $\Omega$ эти величины равны $i_x=i\cos\Omega$ и
$i_y=i\sin\Omega$ в первом порядке. Обратные выражения совпадают
с двумя последними строками \eqref{eq:free-elements-solution}.
При $i=0$ обе компоненты равны нулю.

Нормированные параметры
$\alpha_a,e_x,e_y,i_x,i_y$ постоянны на каждом пассивном участке.
При наличии управления их производные получаются
дифференцированием \eqref{eq:semimajor-state},
\eqref{eq:eccentricity-state} и \eqref{eq:inclination-state}
с подстановкой \eqref{eq:linear-dynamics}:
\begin {equation}
 \boxed{\begin {aligned}
  \dot\alpha_a&=\frac{{2T}}{{V_0}},\\
  \dot e_x&=\frac{{S\sin\varphi+2T\cos\varphi}}{{V_0}},&
  \dot e_y&=\frac{{-S\cos\varphi+2T\sin\varphi}}{{V_0}},\\
  \dot i_x&=\frac{{W\cos\varphi}}{{V_0}},&
  \dot i_y&=\frac{{W\sin\varphi}}{{V_0}}.
 \end {aligned}}
 \label {eq:linear-element-dynamics}
\end {equation}
Например, если $h=\Delta\widehat r+2\Delta\widehat V_t$, то
$\dot h=-n_0\Delta\widehat V_r+2T/V_0$, $\Delta\dot{\widehat V}_r=n_0h+S/V_0$.
В производной $e_x=h\cos\varphi+\Delta\widehat V_r\sin\varphi$ все члены свободного
движения сокращаются, а остаются только составляющие управления.
Так же выводятся остальные строки. В отличие от уравнений
классических элементов, система
\eqref{eq:linear-element-dynamics} не содержит деления на
эксцентриситет или на синус наклонения.

\subsection{Изменение параметров от одного импульса}

Поскольку координаты при импульсе непрерывны, в
\eqref{eq:semimajor-state}, \eqref{eq:eccentricity-state} и
\eqref{eq:inclination-state} изменяются только $\Delta\widehat V_r$, $\Delta\widehat V_t$ и $\widehat V_z$.
Угол точки приложения равен $\varphi_j$. Скачки орбитальных параметров
имеют вид
\begin {equation}
 \boxed{\begin {aligned}
  \delta\alpha_{a,j}&=2{\Delta\widehat V_{tj}},\\
  \delta e_{xj}&={\Delta\widehat V_{rj}}\sin\varphi_j+2{\Delta\widehat V_{tj}}\cos\varphi_j,\\
  \delta e_{yj}&=-{\Delta\widehat V_{rj}}\cos\varphi_j+2{\Delta\widehat V_{tj}}\sin\varphi_j,\\
  \delta i_{xj}&={\Delta\widehat V_{zj}}\cos\varphi_j,\\
  \delta i_{yj}&={\Delta\widehat V_{zj}}\sin\varphi_j.
 \end {aligned}}
 \label {eq:impulse-map}
\end {equation}
Здесь $\delta$ обозначает скачок при отдельном импульсе;
например, $\delta e_{xj}=e_x(t_j+0)-e_x(t_j-0)$.
Те же формулы можно получить интегрированием
\eqref{eq:linear-element-dynamics} по короткому интервалу
тяги, на котором угол практически не меняется.

Из \eqref{eq:impulse-map} следуют полезные частные случаи.
Трансверсальный импульс изменяет большую полуось и создаёт
добавку к вектору эксцентриситета вдоль радиуса точки приложения.
Радиальный импульс в первом порядке не меняет большую полуось и
создаёт добавку к вектору эксцентриситета в направлении
$-\boldsymbol e_t$. Нормальный импульс изменяет плоскость и
не меняет плоские параметры в первом порядке. Например,
при $\varphi_j=0$ положительный радиальный импульс даёт
$\delta e_y=-{\Delta\widehat V_{rj}}$, а положительный нормальный импульс ---
$\delta i_x={\Delta\widehat V_{zj}}$. Эти проверки фиксируют знаки принятой системы
координат.

Отсутствие линейного изменения большой полуоси при радиальном или
нормальном импульсе не означает точного сохранения энергии.
В исходной задаче
\[
 \delta\mathcal E=\boldsymbol V(t_j-0)\mathbin{\cdot}
 \Delta\boldsymbol V_j+\frac{1}{2}\lVert\Delta\boldsymbol V_j\rVert^2.
\]
Второе слагаемое квадратично по импульсу. При вычислении первой
вариации около круговой орбиты первое слагаемое равно
$V_0\Delta V_{tj}$, а все остальные вклады имеют второй порядок.

\subsection{Пять уравнений перехода и свободная конечная фаза}
\label{sec:five-boundaries}

Пусть переход содержит $N$ импульсов. Между ними орбитальные
параметры постоянны, а их скачки определяются
\eqref{eq:impulse-map}. Поэтому суммарное изменение равно сумме
изменений от отдельных импульсов. Обозначим заданные разности
\begin {equation}
 {\Delta a}=\alpha_{a,f}-\alpha_{a,i}=\frac{{a_f-a_i}}{{r_0}},\qquad
 {\Delta e_x}=e_{xf}-e_{xi},\qquad {\Delta e_y}=e_{yf}-e_{yi}.
 \label {eq:boundary-orbit-differences}
\end {equation}
В частности, $\sqrt{{(\Delta e_x)}^2+{(\Delta e_y)}^2}$ является длиной \emph{разности}
векторов эксцентриситета и обычно не равно $|e_f-e_i|$.
При одинаковых скалярных эксцентриситетах орбиты могут иметь разные
направления перицентров и ненулевые ${\Delta e_x},{\Delta e_y}$.

Для пространственного перехода направим неподвижную ось $X$ по
линии пересечения начальной и конечной плоскостей, к восходящему
узлу конечной орбиты относительно начальной. Пусть ${\Delta i}\geq0$ ---
малый угол между плоскостями. Тогда начальная плоскость имеет
$i_{xi}=i_{yi}=0$, а конечная --- $i_{xf}={\Delta i}$, $i_{yf}=0$.
Действительно, её высота и нормальная скорость в первом порядке
равны $z/r_0={\Delta i}\sin\varphi$ и $V_z/V_0={\Delta i}\cos\varphi$.
При совпадающих плоскостях ${\Delta i}=0$ ось $X$ можно выбрать произвольно.
Если исходные элементы заданы относительно экватора или другой
внешней плоскости, сначала следует найти взаимную ориентацию
плоскостей и пересчитать компоненты в общие оси. Разность двух
наклонений к экватору сама по себе не определяет ${\Delta i}$.

Граничные условия перехода принимают вид
\begin {equation}
 \boxed{\begin {aligned}
  \sum_{j=1}^{N}2{\Delta\widehat V_{tj}}&={\Delta a},\\
  \sum_{j=1}^{N}({\Delta\widehat V_{rj}}\sin\varphi_j+2{\Delta\widehat V_{tj}}\cos\varphi_j)&={\Delta e_x},\\
  \sum_{j=1}^{N}(-{\Delta\widehat V_{rj}}\cos\varphi_j+2{\Delta\widehat V_{tj}}\sin\varphi_j)&={\Delta e_y},\\
  \sum_{j=1}^{N}{\Delta\widehat V_{zj}}\sin\varphi_j&=0,\\
  \sum_{j=1}^{N}{\Delta\widehat V_{zj}}\cos\varphi_j&={\Delta i}.
 \end {aligned}}
 \label {eq:five-boundary-equations}
\end {equation}
В матричной записи
\begin {equation}
 \sum_{j=1}^{N} B(\varphi_j)\Delta\widehat{\boldsymbol V}_j=\boldsymbol b,\qquad
 B(\varphi)=
 \begin {pmatrix}
  0&2&0\\
  \sin\varphi&2\cos\varphi&0\\
  -\cos\varphi&2\sin\varphi&0\\
  0&0&\sin\varphi\\
  0&0&\cos\varphi
 \end {pmatrix},\qquad
 \boldsymbol b=\begin {pmatrix}{\Delta a}\\{\Delta e_x}\\{\Delta e_y}\\0\\{\Delta i}\end {pmatrix}.
 \label {eq:boundary-map}
\end {equation}
Эта нормировка и состав граничных параметров соответствуют
используемым Барановым условиям перехода
\cite{baranov-ru,baranov-en}; отсчёт всех углов здесь единый и
задан явно.

Почему при шестимерном состоянии получено пять условий?
Пять параметров задают орбиту как геометрическую кривую и её
масштаб, но не положение аппарата на ней. Введём фазовую переменную
\begin {equation}
 \ell=\Delta u-2e_x\sin\varphi+2e_y\cos\varphi.
 \label {eq:mean-phase-variable}
\end {equation}
В первом порядке она равна отклонению средней долготы аппарата от
фазы опорного круга. Подставляя
\eqref{eq:linear-dynamics} и \eqref{eq:linear-element-dynamics},
получаем
\begin {equation}
 \dot\ell=-\frac{3}{2}n_0\alpha_a-\frac{{2S}}{{V_0}},\qquad
 \delta\ell_j=-2{\Delta\widehat V_{rj}}.
 \label {eq:phase-dynamics}
\end {equation}
Например, на пассивном участке первое равенство согласуется с
линеаризацией зависимости среднего движения от большой полуоси:
$n(a)-n_0=-(3/2)n_0\alpha_a$. При импульсе физическая долгота
не меняется, но средняя долгота может измениться, поскольку
изменяются параметры оскулирующей орбиты.

Если бы требовалось встретить заданный аппарат в назначенный момент
$t_f$, пришлось бы дополнительно выполнить условие на фазу.
При импульсном управлении интегрирование
\eqref{eq:phase-dynamics} от $t_i$ до $t_f$ даёт
\begin {equation}
 \ell_f-\ell_i=-\frac{3}{2}n_0\alpha_{a,i}(t_f-t_i)
 -\sum_{j=1}^{N}\bigl[2{\Delta\widehat V_{rj}}+3n_0{\Delta\widehat V_{tj}}(t_f-t_j)\bigr].
 \label {eq:phase-boundary-example}
\end {equation}
В задаче перехода на орбиту конечная фаза не назначается; свободна
также начальная точка приложения манёвра на исходной орбите.
Поэтому \eqref{eq:phase-boundary-example} определяет получившуюся
фазу, но не является ограничением на управление. Исключение
шестого условия связано с постановкой задачи и не означает
отсутствия продольного движения.

Итак, в линейной модели требуется выбрать число импульсов, их
углы и компоненты так, чтобы выполнить
\eqref{eq:boundary-map} и минимизировать
\begin {equation}
 J=\sum_{j=1}^{N}{\Delta\widehat V_j},\qquad
 {\Delta\widehat V_j}=\lVert\Delta\widehat{\boldsymbol V}_j\rVert
 =\sqrt{{(\Delta\widehat V_{rj})}^2+{(\Delta\widehat V_{tj})}^2+{(\Delta\widehat V_{zj})}^2},\qquad
 \Delta V_\Sigma=V_0J.
 \label {eq:linear-transfer-cost}
\end {equation}
Величина ${\Delta\widehat V_j}$ неотрицательна; знак имеют отдельные компоненты
импульса. В компланарной задаче ${\Delta i}=0$ и нормальные компоненты
можно положить равными нулю, оставив первые три уравнения.

\subsection{Границы применимости полученной модели}

Уравнения \eqref{eq:linear-dynamics} описывают малые отклонения
от конкретного опорного движения. Формулы
\eqref{eq:impulse-map} и \eqref{eq:boundary-map} описывают
первый порядок изменений орбиты около того же круга. Для их
применения недостаточно только малости эксцентриситетов:
необходимы также малость $|a-r_0|/r_0$, взаимного наклонения
и нормированных импульсов на всех участках перехода. Например,
две круговые орбиты с сильно различающимися радиусами имеют
$e_i=e_f=0$, но не принадлежат общей малой окрестности выбранного
круга.

Линеаризация также не даёт равномерной точности на произвольно
больших временах. Член $-(3/2)\alpha_a(\varphi-\varphi_*)$ в
\eqref{eq:free-phase} накапливает фазовое расхождение.
Если $\varepsilon n_0|t-t_*|$ перестаёт быть малым, приближение
координат относительно одной опорной точки требует пересмотра.
В задаче со свободными фазами матрицу $B(\varphi_j)$ можно применять
локально в выбранных геометрических точках: замена фактического
угла на угол соответствующей точки опорного круга изменяет
первый порядок импульсного эффекта только на величину второго
порядка, пока угловая ошибка мала. Добавление произвольного
числа витков к времени полёта не сопровождается автоматически
гарантией малой ошибки полного линейного прогноза.

Наконец, найденные по пяти линейным условиям импульсы обеспечивают
выход на заданную орбиту лишь в первом приближении. Проверка в
точной модели должна заново вычислять оскулирующие элементы после
последнего импульса; при необходимости параметры манёвров
уточняются по фактической невязке. Доказательство минимальности
$J$ в модели \eqref{eq:boundary-map} относится именно к этой
модели. Этот предел вывода будет существенен при формулировке
условий оптимальности.

% ==================== optimality ====================
\section{Условия оптимальности}
\label{sec:optimality}

В предыдущем разделе получены линейные уравнения перехода. Теперь
необходимо выбрать число импульсов, их положения и векторы так, чтобы
минимизировать суммарную характеристическую скорость. Сначала выведем
условия оптимальности непосредственно для этой задачи. Затем установим
их связь с сопряжёнными уравнениями Лоудена для точного движения в
центральном поле \cite[гл.~3, \S~5.1]{lawden}.

\subsection{Задача с произвольным числом импульсов}

В безразмерных переменных требуется найти
\begin {equation}
 \begin {gathered}
 J_* = \min_{N,\,\varphi_j,\,\Delta\widehat{\boldsymbol V}_j}
       \sum_{j=1}^{N}\lVert\Delta\widehat{\boldsymbol V}_j\rVert,\\
 \sum_{j=1}^{N}B(\varphi_j)\Delta\widehat{\boldsymbol V}_j=\boldsymbol b,
 \qquad \boldsymbol b=({\Delta a},{\Delta e_x},{\Delta e_y},0,{\Delta i})^{\mathsf T}.
 \end {gathered}
 \label {eq:linear-optimization}
\end {equation}
Здесь $B(\varphi)$ --- матрица из уравнений перехода,
$\varphi_j$ определены по модулю $2\pi$,
$\Delta\widehat{\boldsymbol V}_j\in\R^3$, а число $N$ заранее не задано.
Нулевой импульс можно удалить, поэтому при обсуждении отдельного
манёвра будем считать ${\Delta\widehat V_j}=\lVert\Delta\widehat{\boldsymbol V}_j\rVert>0$.
При $\boldsymbol b=0$ решение очевидно: $J_*=0$ и импульсы отсутствуют.
Далее рассматривается $\boldsymbol b\ne0$.

Свобода выбора фаз в \eqref{eq:linear-optimization} связана с принятой
постановкой: время перехода и конечная фаза не ограничены. В линейной
модели импульсы с заданными фазами можно расположить в хронологическом
порядке, при необходимости добавляя обороты. Этот факт не означает,
что приближение остаётся равномерно точным при произвольно большом
времени движения. Погрешность восстановления реальной траектории
оценивается отдельно.

Представим ненулевой импульс в виде
\begin {equation}
 \Delta\widehat{\boldsymbol V}_j={\Delta\widehat V_j}\boldsymbol d_j,
 \qquad {\Delta\widehat V_j}>0,\qquad \lVert\boldsymbol d_j\rVert=1.
\end {equation}
Вектор $B(\varphi_j)\boldsymbol d_j\in\R^5$ показывает, какие изменения
параметров орбиты даёт единичная затрата характеристической скорости
в выбранной точке и направлении. Введём множество таких изменений
\begin {equation}
 \mathcal S=
 \bigl\{B(\varphi)\boldsymbol d:
          0\leq\varphi\leq2\pi,\ \lVert\boldsymbol d\rVert=1\bigr\}
\end {equation}
и его выпуклую оболочку $\mathcal C=\operatorname {conv}\mathcal S$.
Под выпуклой оболочкой понимается множество всех конечных средних
$\sum_j w_j\boldsymbol s_j$, где
$\boldsymbol s_j\in\mathcal S$, $w_j\geq0$ и $\sum_jw_j=1$.
Если суммарная затрата равна $J$, то
\begin {equation}
 \frac{{\boldsymbol b}}{{J}}
   =\sum_j\frac{{{\Delta\widehat V_j}}}{{J}}B(\varphi_j)\boldsymbol d_j
   \in\mathcal C.
 \label {eq:reachable-convex}
\end {equation}
Обратно, любое такое среднее задаёт набор импульсов
$\Delta\widehat{\boldsymbol V}_j=Jw_j\boldsymbol d_j$ с затратой $J$.
Следовательно, поиск оптимального управления эквивалентен поиску
наименьшего множителя $J$, для которого $\boldsymbol b\in J\mathcal C$.

Множество $\mathcal C$ компактно, симметрично относительно начала
координат и содержит начало координат внутри. Последнее свойство
существенно: все пять требуемых приращений доступны при конечной
затрате. Его можно проверить без анализа отдельных манёвров.
Если вектор $\boldsymbol l\in\R^5$ ортогонален всем элементам
$\mathcal S$, то $B(\varphi)^{\mathsf T}\boldsymbol l=0$ при любом $\varphi$.
Из первой и второй компонент этого тождества следуют
$l_x=l_y=l_a=0$, а из третьей --- $l_s=l_c=0$.
Таким образом, $\mathcal S$ порождает всё пространство $\R^5$;
его симметричная выпуклая оболочка имеет непустую внутренность.
Компактность можно пояснить следующим образом. Любое выпуклое
сочетание более чем шести точек в $\R^5$ допускает сокращение числа
слагаемых за счёт их аффинной зависимости. Повторяя сокращение,
получаем не более шести точек. Множество всех таких сочетаний
компактно, поскольку $\mathcal S$ компактно.
Число шесть здесь является общей геометрической оценкой; специальная
структура орбитального перехода позволяет уменьшить его.

Поэтому минимум $J_*>0$ существует, а точка
$\boldsymbol b/J_*$ лежит на границе $\mathcal C$.
Через граничную точку выпуклого тела можно провести опорную
гиперплоскость: всё тело находится по одну сторону от неё.
В рассматриваемом случае это означает существование
$\boldsymbol l\ne0$, для которого после выбора масштаба
\begin {equation}
 \boldsymbol l^{\mathsf T}\boldsymbol s\leq1
       \quad(\boldsymbol s\in\mathcal C),
 \qquad
 \boldsymbol l^{\mathsf T}\boldsymbol b=J_*.
 \label {eq:supporting-plane}
\end {equation}
Положительность нормирующей константы следует из того, что начало
координат находится внутри $\mathcal C$.
Геометрический факт об опорной плоскости можно получить, приближая
граничную точку внешними точками. Для каждой внешней точки выбирают
ближайшую точку выпуклого тела. Вектор между ними образует неострый
угол с любым направлением из ближайшей точки внутрь тела; иначе
расстояние можно было бы уменьшить. Предельное положение
перпендикулярной ему плоскости и даёт опорную плоскость.

\subsection{Базис-вектор и необходимое условие для всех допустимых фаз}

Множители $l_a,l_x,l_y,l_s,l_c$ задают линейную оценку изменений
параметров орбиты. Если единичное направление импульса равно
$\boldsymbol d$, импульс величины $\varepsilon$ создаёт изменение
$\delta\boldsymbol b=\varepsilon B(\varphi)\boldsymbol d$.
Его оценка равна
\[
 \boldsymbol l^{\mathsf T}\delta\boldsymbol b
 =\varepsilon\bigl(B(\varphi)^{\mathsf T}\boldsymbol l\bigr)^{\mathsf T}
 \boldsymbol d.
\]
Таким образом, множители граничных условий можно перенести из
пространства орбитальных параметров в пространство трёх компонент
импульса. Полученный вектор и задаёт оценку эффективности его направления.

Определим \emph{базис-вектор}
\begin {equation}
 \boldsymbol p(\varphi)=B(\varphi)^{\mathsf T}\boldsymbol l,
 \qquad
 \boldsymbol l=(l_a,l_x,l_y,l_s,l_c)^{\mathsf T}.
 \label {eq:primer-definition}
\end {equation}
В принятой нормировке все его компоненты безразмерны:
\begin {equation}
 \boxed{\begin {aligned}
 p_r(\varphi)&=l_x\sin\varphi-l_y\cos\varphi,\\
 p_t(\varphi)&=2l_a+2l_x\cos\varphi+2l_y\sin\varphi,\\
 p_z(\varphi)&=l_s\sin\varphi+l_c\cos\varphi.
 \end {aligned}}
 \label {eq:primer-components}
\end {equation}
Пять коэффициентов $\boldsymbol l$ постоянны. Это связано с тем, что
в выбранных орбитальных переменных линейная матрица действия импульса
зависит от фазы, но не от уже накопленных приращений орбиты.

Первое неравенство \eqref{eq:supporting-plane}, применённое к
$\boldsymbol s=B(\varphi)\boldsymbol d$, принимает вид
\begin {equation}
 \boldsymbol p(\varphi)^{\mathsf T}\boldsymbol d\leq1
        \qquad\text {при любом }\lVert\boldsymbol d\rVert=1.
\end {equation}
Наибольшее скалярное произведение с единичным вектором равно модулю
самого вектора. Следовательно,
\begin {equation}
 \lVert\boldsymbol p(\varphi)\rVert\leq1
       \qquad\text {при всех }\varphi.
 \label {eq:primer-global-bound}
\end {equation}
В оптимальном разложении \eqref{eq:reachable-convex} среднее
скалярных произведений равно единице. Поскольку каждое из них не
превышает единицы, любое слагаемое с положительным весом также
должно давать единицу. Равенство в неравенстве Коши достигается
только при совпадении направлений. Поэтому для каждого ненулевого
оптимального импульса
\begin {equation}
 \lVert\boldsymbol p(\varphi_j)\rVert=1,
 \qquad
 \Delta\widehat{\boldsymbol V}_j={\Delta\widehat V_j}\boldsymbol p(\varphi_j),
 \qquad {\Delta\widehat V_j}>0.
 \label {eq:primer-contact}
\end {equation}
Таким образом, импульсы располагаются в точках абсолютного максимума
модуля базис-вектора и направлены вдоль него. Все используемые
максимумы имеют одну высоту, равную единице.
Условие \eqref{eq:primer-contact} задаёт \emph{направление}, но не
величину импульса. Величины ${\Delta\widehat V_j}$ определяются из граничных уравнений.

Вывод через выпуклую оболочку учитывает произвольное число импульсов.
Это устраняет существенный пробел, который возник бы при
дифференцировании задачи с заранее заданным числом манёвров:
стационарность внутри такой схемы сама по себе не исключает её
улучшения добавлением нового импульса.

\subsection{Первая вариация по величинам, направлениям и положениям импульсов}

Условия в используемых точках можно дополнительно получить
непосредственным варьированием. Для фиксированного числа ненулевых
импульсов составим функцию Лагранжа
\begin {equation}
 \mathcal L=
 \sum_j\lVert\Delta\widehat{\boldsymbol V}_j\rVert
 -\boldsymbol l^{\mathsf T}
       \left(\sum_j B(\varphi_j)\Delta\widehat{\boldsymbol V}_j-\boldsymbol b\right).
\end {equation}
Поскольку
$\delta\lVert\Delta\widehat{\boldsymbol V}_j\rVert
  =\Delta\widehat{\boldsymbol V}_j^{\mathsf T}\delta\Delta\widehat{\boldsymbol V}_j/{\Delta\widehat V_j}$,
первая вариация имеет вид
\begin {equation}
 \delta\mathcal L=
 \sum_j\left(\frac{{\Delta\widehat{\boldsymbol V}_j}}{{{\Delta\widehat V_j}}}
             -\boldsymbol p(\varphi_j)\right)^{\mathsf T}
                \delta\Delta\widehat{\boldsymbol V}_j
 -\sum_j \boldsymbol p'(\varphi_j)^{\mathsf T}
                  \Delta\widehat{\boldsymbol V}_j\,\delta\varphi_j,
 \label {eq:impulse-first-variation}
\end {equation}
где штрих означает производную по $\varphi$.
Произвольность вариаций трёх компонент $\Delta\widehat{\boldsymbol V}_j$ даёт
\eqref{eq:primer-contact}; произвольность вариации внутренней фазы
даёт
\begin {equation}
 \boldsymbol p(\varphi_j)^{\mathsf T}\boldsymbol p'(\varphi_j)=0,
 \qquad
 \left.\frac{{\dd}}{{\dd\varphi}}
     \lVert\boldsymbol p(\varphi)\rVert\right|_{\varphi=\varphi_j}=0.
 \label {eq:primer-stationarity}
\end {equation}
Производная самого вектора при этом, вообще говоря, не равна нулю.
Базис-вектор может поворачиваться, сохраняя свой модуль.
Кроме того, стационарная точка может оказаться минимумом либо
не самым высоким максимумом. Поэтому
\eqref{eq:primer-stationarity} не заменяет глобальную проверку
\eqref{eq:primer-global-bound}.

Роль дополнительного импульса видна из его вклада в функцию Лагранжа.
При добавлении $\varepsilon\boldsymbol d$ в фазе $\varphi$, где
$\varepsilon>0$ и $\lVert\boldsymbol d\rVert=1$, вклад равен
\begin {equation}
 \varepsilon\bigl(1-
       \boldsymbol p(\varphi)^{\mathsf T}\boldsymbol d\bigr).
\end {equation}
Наиболее выгодное направление совпадает с $\boldsymbol p$; тогда
коэффициент при $\varepsilon$ равен
$1-\lVert\boldsymbol p(\varphi)\rVert$.
В данном выводе общий множитель $\boldsymbol l$ и невозможность
отрицательного коэффициента уже обоснованы
\eqref{eq:supporting-plane}. Одна только подстановка нового импульса
в исходное управление была бы недостаточна: она нарушает конечные
условия, которые также необходимо учитывать.

Если импульсы допускаются лишь на заданном временном интервале,
в его фиксированных концах двусторонняя вариация времени отсутствует.
Тогда равенство \eqref{eq:primer-stationarity} для концевого импульса
не обязательно: возникают односторонние условия.
В нашей задаче разрешена любая фаза на окружности, поэтому такого
конца допустимого интервала фаз нет.

\subsection{Достаточность и нижняя оценка характеристической скорости}

\begin{theorem}
Для линейной задачи \eqref{eq:linear-optimization} с
$\boldsymbol b\ne0$ набор импульсов глобально оптимален тогда и
только тогда, когда существуют постоянные $\boldsymbol l$,
для которых выполнены условия
\begin {equation}
 \boxed{\begin {gathered}
   \sum_jB(\varphi_j)\Delta\widehat{\boldsymbol V}_j=\boldsymbol b,\\
   \boldsymbol p(\varphi)=B(\varphi)^{\mathsf T}\boldsymbol l,
   \qquad \lVert\boldsymbol p(\varphi)\rVert\leq1
                              \quad\text {при всех }\varphi,\\
   \Delta\widehat{\boldsymbol V}_j={\Delta\widehat V_j}\boldsymbol p(\varphi_j),
   \qquad {\Delta\widehat V_j}>0,\qquad
   \lVert\boldsymbol p(\varphi_j)\rVert=1.
 \end {gathered}}
 \label {eq:certificate}
\end {equation}
\end{theorem}

\begin{proof}
Необходимость установлена построением опорной гиперплоскости.
Для доказательства достаточности рассмотрим любой другой
допустимый набор $\{\varphi_k^{\mathrm {alt}},\Delta\widehat{\boldsymbol V}^{\mathrm {alt}}_k\}$,
в том числе с другим числом импульсов. Имеем
\begin {align}
 J^{\mathrm {alt}}
 &=\sum_k\lVert\Delta\widehat{\boldsymbol V}^{\mathrm {alt}}_k\rVert
 \geq\sum_k\boldsymbol p(\varphi_k^{\mathrm {alt}})^{\mathsf T}
                      \Delta\widehat{\boldsymbol V}^{\mathrm {alt}}_k \notag\\
 &=\boldsymbol l^{\mathsf T}
         \sum_k B(\varphi_k^{\mathrm {alt}})\Delta\widehat{\boldsymbol V}^{\mathrm {alt}}_k
  =\boldsymbol l^{\mathsf T}\boldsymbol b.
 \label {eq:dual-lower-bound}
\end {align}
Для проверяемого набора все неравенства обращаются в равенства:
\begin {equation}
 \boldsymbol l^{\mathsf T}\boldsymbol b
   =\sum_j{\Delta\widehat V_j}\lVert\boldsymbol p(\varphi_j)\rVert^2
   =\sum_j{\Delta\widehat V_j}=J.
\end {equation}
Следовательно, $J^{\mathrm {alt}}\geq J$ для любого допустимого управления.
\end{proof}

Это доказательство использует линейность граничных уравнений и
выпуклость нормы. Для точного нелинейного перехода те же условия
на базис-вектор в общем случае являются лишь необходимыми:
линейная оценка \eqref{eq:dual-lower-bound} уже не связывает
произвольные конечные орбиты с одним и тем же оператором $B$.

Полезна также двойственная запись оптимальной цены:
\begin {equation}
 J_*(\boldsymbol b)=
 \max_{\boldsymbol l}
 \left\{\boldsymbol l^{\mathsf T}\boldsymbol b:
       \max_\varphi\lVert B(\varphi)^{\mathsf T}\boldsymbol l\rVert
           \leq1\right\}.
 \label {eq:dual-problem}
\end {equation}
Если $J_*$ дифференцируема по всем пяти компонентам
$\boldsymbol b$, то оптимальный множитель равен её градиенту:
\begin {equation}
 \boldsymbol l=\nabla_{\boldsymbol b}J_*.
 \label {eq:cost-gradient}
\end {equation}
Действительно, тот же допустимый множитель даёт для изменённого
задания нижнюю оценку
$J_*(\boldsymbol b+\delta\boldsymbol b)
 \geq J_*(\boldsymbol b)+\boldsymbol l^{\mathsf T}\delta\boldsymbol b$.
В точке дифференцируемости такая линейная опорная функция имеет
коэффициенты градиента. На границе ветвей множитель может быть
неединственным; тогда использовать производную несуществующей
гладкой ветви нельзя.

Нулевое значение четвёртой компоненты задания не означает
$l_s=0$. Производная в \eqref{eq:cost-gradient} берётся по
полному вектору изменений, включая обе компоненты наклонения.
Если до дифференцирования повернуть оси вдоль линии узлов,
при вариации направления конечной плоскости необходимо учесть
и изменение этого поворота.

\subsection{Связь с сопряжёнными уравнениями Лоудена}

Покажем, откуда появляется та же векторная функция в общей задаче.
В этом подразделе $\boldsymbol R$ и $\boldsymbol V$ --- размерные
инерциальные векторы, $R=\lVert\boldsymbol R\rVert$, а $\Delta\boldsymbol V_j$ --- размерные
импульсы. На баллистических дугах
\begin {equation}
 \dot{\boldsymbol R}=\boldsymbol V,
 \qquad \dot{\boldsymbol V}=\boldsymbol g(\boldsymbol R),
 \qquad
 \boldsymbol g(\boldsymbol R)=-\frac{{\mu}}{{R^3}}\boldsymbol R.
\end {equation}
Минимизируется $\sum_j\lVert\Delta\boldsymbol V_j\rVert$.
Для регулярной, или нормальной, экстремали коэффициент при этом
функционале в расширенной вариационной задаче можно принять равным
единице. Нормальность означает, что множители не описывают только
вырожденные ограничения при нулевом коэффициенте стоимости.
Для линейной задачи необходимость условий уже доказана выше
без дополнительного предположения о нормальности.

Введём сопряжённые векторы $\boldsymbol a(t)$ и
$\boldsymbol p(t)$ и добавим к стоимости интегралы по всем
баллистическим дугам:
\begin {equation}
 \mathcal A=\sum_j\lVert\Delta\boldsymbol V_j\rVert
 +\sum_{\text{дуги}}\int
 \left[
 \boldsymbol a^{\mathsf T}(\dot{\boldsymbol R}-\boldsymbol V)
 +\boldsymbol p^{\mathsf T}
          (\dot{\boldsymbol V}-\boldsymbol g(\boldsymbol R))
 \right]\dd t.
 \label {eq:augmented-coast-action}
\end {equation}
Множители концевых ограничений здесь не выписаны, поскольку сначала
рассматриваются внутренние вариации. Варьирование $\boldsymbol R$
и $\boldsymbol V$ и интегрирование слагаемых с производными по
частям дают внутри дуг
\begin {equation}
 \dot{\boldsymbol a}
       =-\left(\frac{{\partial\boldsymbol g}}{{\partial\boldsymbol R}}\right)^{\mathsf T}
          \boldsymbol p,
 \qquad
 \dot{\boldsymbol p}=-\boldsymbol a.
\end {equation}
Исключая $\boldsymbol a$, получаем уравнение базис-вектора:
\begin {equation}
 \ddot{\boldsymbol p}
  =\left(\frac{{\partial\boldsymbol g}}{{\partial\boldsymbol R}}\right)^{\mathsf T}
       \boldsymbol p
  =-\frac{{\mu}}{{R^3}}\boldsymbol p
    +\frac{{3\mu}}{{R^5}}
           (\boldsymbol p^{\mathsf T}\boldsymbol R)\boldsymbol R.
 \label {eq:lawden-primer-ode}
\end {equation}
Второе равенство следует из покомпонентного дифференцирования
$g_i=-\mu R_i/R^3$:
$\partial g_i/\partial R_k
 =-\mu\delta_{ik}/R^3+3\mu R_i R_k/R^5$.
Матрица производных симметрична.
Уравнение \eqref{eq:lawden-primer-ode} линейно по
$\boldsymbol p$ вдоль заданной баллистической дуги. Сама задача
поиска этой дуги между двумя орбитами от этого линейной не становится.

В момент импульса
$\boldsymbol R^+=\boldsymbol R^-$ и
$\boldsymbol V^+=\boldsymbol V^-+\Delta\boldsymbol V$.
Вариации положения и доимпульсной скорости в граничных членах
\eqref{eq:augmented-coast-action} входят в выражение
\begin {equation}
 \begin {aligned}
 \delta\mathcal A_j={}&(\boldsymbol a^--\boldsymbol a^+)^{\mathsf T}
 \delta\boldsymbol R
 +(\boldsymbol p^--\boldsymbol p^+)^{\mathsf T}\delta\boldsymbol V^-\\
 &+\left(\frac{{\Delta\boldsymbol V}}{{\lVert\Delta\boldsymbol V\rVert}}
 -\boldsymbol p^+\right)^{\mathsf T}\delta\Delta\boldsymbol V.
 \end {aligned}
\end {equation}
Здесь использовано равенство
$\delta\boldsymbol V^+=\delta\boldsymbol V^-+\delta\Delta\boldsymbol V$.
Независимость трёх указанных вариаций даёт
$\boldsymbol a^+=\boldsymbol a^-$ и
$\boldsymbol p^+=\boldsymbol p^-$.
Коэффициент при вариации самого импульса равен
\begin {equation}
 \frac{{\Delta\boldsymbol V}}{{\lVert\Delta\boldsymbol V\rVert}}
          -\boldsymbol p.
\end {equation}
Следовательно, $\boldsymbol p$ направлен вдоль ненулевого импульса
и имеет в этой точке единичный модуль. Поскольку
$\dot{\boldsymbol p}=-\boldsymbol a$, первая инерциальная производная
базис-вектора также непрерывна. Утверждать непрерывность всех его
производных оснований нет.

Для свободного внутреннего момента импульса возникает условие
непрерывности гамильтониана
\begin {equation}
 H=\boldsymbol a^{\mathsf T}\boldsymbol V
            +\boldsymbol p^{\mathsf T}\boldsymbol g.
\end {equation}
Оно следует из коэффициента при перемещении границы двух
баллистических интегралов. Так как положение и сопряжённые векторы
непрерывны,
\begin {equation}
 0=H^+-H^-
  =\boldsymbol a^{\mathsf T}\Delta\boldsymbol V
  =-\dot{\boldsymbol p}^{\mathsf T}\Delta\boldsymbol V.
\end {equation}
После подстановки направления импульса получаем
$\boldsymbol p^{\mathsf T}\dot{\boldsymbol p}=0$, то есть
стационарность его модуля. Это размерный инерциальный аналог
\eqref{eq:primer-stationarity}.

Наконец, для малого дополнительного импульса точная первая вариация
расширенного функционала для нормальной экстремали содержит
$\lVert\delta\boldsymbol V\rVert-
\boldsymbol p^{\mathsf T}\delta\boldsymbol V$.
Её неотрицательность при произвольном направлении даёт
$\lVert\boldsymbol p\rVert\leq1$.
Таким образом, условия Лоудена согласуются с полученными выше
условиями, но их достаточность для нелинейной задачи из этой первой
вариации не следует.

\subsection{Базис-вектор на опорной круговой орбите}

Положим теперь $R=r_0$ и рассмотрим вращающиеся единичные векторы
$\boldsymbol e_r$, $\boldsymbol e_t$, $\boldsymbol e_z$,
для которых
\begin {equation}
 \dot{\boldsymbol e}_r=n_0\boldsymbol e_t,
 \qquad \dot{\boldsymbol e}_t=-n_0\boldsymbol e_r,
 \qquad \dot{\boldsymbol e}_z=0.
\end {equation}
Компоненты $p_r,p_t,p_z$ определены разложением
$\boldsymbol p=p_r\boldsymbol e_r+p_t\boldsymbol e_t+p_z\boldsymbol e_z$.
При дифференцировании необходимо учитывать вращение осей:
\begin {equation}
 \dot{\boldsymbol p}
 = (\dot p_r-n_0p_t)\boldsymbol e_r
  +(\dot p_t+n_0p_r)\boldsymbol e_t
  +\dot p_z\boldsymbol e_z.
\end {equation}
Подстановка второй производной в
\eqref{eq:lawden-primer-ode} и переход от времени к
$\varphi=n_0t+\varphi_0$ дают
\begin {equation}
 p_r''-2p_t'-3p_r=0,
 \qquad p_t''+2p_r'=0,
 \qquad p_z''+p_z=0.
 \label {eq:circular-primer-ode}
\end {equation}
Из второго уравнения $p_t'+2p_r=C$.
Тогда первое превращается в $p_r''+p_r=2C$.
Общее решение имеет шесть постоянных:
\begin {equation}
 \begin {aligned}
 p_r&=c_1\cos\varphi+c_2\sin\varphi+2C,\\
 p_t&=2c_2\cos\varphi-2c_1\sin\varphi-3C\varphi+c_3,\\
 p_z&=c_4\cos\varphi+c_5\sin\varphi.
 \end {aligned}
 \label {eq:circular-primer-general}
\end {equation}

Появление секулярного члена $-3C\varphi$ не противоречит
\eqref{eq:primer-components}: ещё не использовано условие
трансверсальности по свободному времени.
В автономной задаче с конечным множеством, не зависящим явно от
времени, и свободным конечным временем это условие есть $H=0$.
Поясним появление этого равенства. Пусть
$\boldsymbol Z=(\boldsymbol R,\boldsymbol V)$ и
$\boldsymbol \Lambda=(\boldsymbol a,\boldsymbol p)$.
Полная вариация конечного состояния связана с вариацией при
фиксированном времени равенством
$\delta\boldsymbol Z_f=\delta\boldsymbol Z(t_f)+
\dot{\boldsymbol Z}(t_f)\delta t_f$.
Поэтому конечный член после интегрирования по частям имеет вид
\[
 \boldsymbol \Lambda_f^{\mathsf T}\delta\boldsymbol Z_f-H_f\delta t_f.
\]
Трансверсальность по конечному состоянию означает, что
$\boldsymbol \Lambda_f$ нормален к допустимому конечному многообразию.
Его скалярное произведение с любой касательной вариацией
$\delta\boldsymbol Z_f$ равно нулю. Так как само многообразие
от времени не зависит, $\delta t_f$ независимо и произвольно:
остаётся $H_f=0$. На пассивных дугах автономной задачи $H$ постоянно,
а при свободных внутренних импульсах непрерывно; следовательно, $H=0$
на всей экстремали. В нашей задаче конечное многообразие является
заданной орбитой со свободной фазой. Её инвариантность относительно
пассивного движения также согласуется с условием $H=0$.
На круговой орбите $\boldsymbol V=V_0\boldsymbol e_t$,
$\boldsymbol g=-n_0V_0\boldsymbol e_r$ и
$\boldsymbol a=-\dot{\boldsymbol p}$, поэтому
\begin {equation}
 H=-n_0V_0(p_t'+2p_r)=-n_0V_0 C.
\end {equation}
Следовательно, $C=0$. Сравнение
\eqref{eq:circular-primer-general} с
\eqref{eq:primer-components} даёт
\begin {equation}
 c_1=-l_y,\qquad c_2=l_x,\qquad c_3=2l_a,
 \qquad c_4=l_c,\qquad c_5=l_s.
\end {equation}
Итак, пятиконстантная форма базис-вектора следует и из точных
сопряжённых уравнений после перехода к опорному кругу и использования
свободного времени. Круговая орбита сама по себе не требует $C=0$.
Для встречи с заданной фазой в заданное время условия
трансверсальности и число граничных ограничений следует выводить
заново.

\subsection{Геометрия годографа и типы оптимальных переходов}
\label{subsec:primer-geometry}

Вначале предположим
$R=\sqrt{l_x^2+l_y^2}>0$ и положим
\begin {equation}
 \begin {gathered}
 \alpha=\operatorname {atan2}(l_y,l_x),
 \qquad \psi=\varphi-\alpha,\qquad D=2l_a,\\
 E=l_s\sin\alpha+l_c\cos\alpha,
 \qquad F=l_s\cos\alpha-l_c\sin\alpha.
 \end {gathered}
\end {equation}
Тогда
\begin {equation}
 \boldsymbol p(\psi)=
 \begin {pmatrix}R\sin\psi\\D+2R\cos\psi\\E\cos\psi+F\sin\psi\end {pmatrix}
 =\begin {pmatrix}0\\D\\0\end {pmatrix}
 +\begin {pmatrix}0\\2R\\E\end {pmatrix}\cos\psi
 +\begin {pmatrix}R\\0\\F\end {pmatrix}\sin\psi.
 \label {eq:rotated-primer}
\end {equation}
Годограф в пространстве трёх \emph{вращающихся компонент} является
эллипсом с центром $(0,D,0)$. Два вектора при синусе и косинусе
задают его плоскость. Его радиально-трансверсальная проекция имеет
полуоси $R$ и $2R$.
Это вспомогательное пространство компонент, а не физическая
траектория конца вектора в неподвижных координатах.

Квадрат модуля равен
\begin {equation}
 \begin {aligned}
 f(\psi)=\lVert\boldsymbol p(\psi)\rVert^2
   &=D^2+R^2+F^2
    +(3R^2+E^2-F^2)\cos^2\psi\\
   &\quad+2EF\sin\psi\cos\psi+4DR\cos\psi.
 \end {aligned}
 \label {eq:primer-norm-quadratic}
\end {equation}
Импульсы возможны только в точках касания этого эллипса с единичной
сферой, причём весь эллипс должен лежать внутри сферы.
Положения и число точек касания определяют классификацию,
используемую Эдельбаумом и Барановым
\cite{edelbaum,baranov-ru,baranov-en}.

\paragraph{Центральный эллипс: узловой случай (NC).}
Если $D=0$, то
$\boldsymbol p(\psi+\pi)=-\boldsymbol p(\psi)$.
У некругового эллипса наиболее удалены от начала координат концы
его большой оси. Таким образом, возможны два противоположных
максимума; их фазы отличаются на $\pi$, а все локальные компоненты
направлений импульсов противоположны.
Для вычисления максимума достаточно максимального собственного
значения симметричной матрицы
\begin {equation}
 Q=\begin {pmatrix}
       4R^2+E^2&EF\\ EF&R^2+F^2
    \end {pmatrix},
 \qquad
 f=(\cos\psi,\sin\psi)Q
        \begin {pmatrix}\cos\psi\\\sin\psi\end {pmatrix}.
\end {equation}
Если оно равно единице, эллипс касается сферы концами большой оси.
Название ``узловой'' будет обосновано при подстановке граничных
условий пространственного перехода: для ${\Delta i}>0$ используемые точки
окажутся на линии взаимных узлов.

\paragraph{Смещённый эллипс: невырожденный случай (NDC).}
Пусть $D\ne0$ и имеются две различные точки абсолютного максимума.
Покажем, почему они должны быть симметричны относительно $\psi=0$.
Функция $1-f(\psi)$ --- неотрицательный тригонометрический многочлен
степени два. В каждой точке касания со сферой он имеет двойной нуль.
Если два нуля обозначить $m-\delta$ и $m+\delta$, такой многочлен
представляется в виде
\begin {equation}
 1-f(\psi)=K[\cos(\psi-m)-\cos\delta]^2,
 \qquad K>0.
 \label {eq:contact-factorization}
\end {equation}
Действительно, произведение
$\sin^2((\psi-m+\delta)/2)
 \sin^2((\psi-m-\delta)/2)$ имеет именно эти два двойных нуля;
раскрытие произведения даёт правую часть с постоянным множителем.
Многочлен степени два с такими нулями определяется с точностью до
этого множителя.
Первая гармоника в \eqref{eq:primer-norm-quadratic} содержит только
$\cos\psi$ и имеет ненулевой коэффициент $4DR$.
Сравнение с \eqref{eq:contact-factorization} даёт
$m=0$ по модулю $\pi$. Следовательно, точки можно записать как
$\psi=\pm\psi_*$. Коэффициент при $\sin2\psi$ должен обращаться
в нуль, то есть $EF=0$.

Если бы $F=0$, зависимость от $c=\cos\psi$ была бы выпуклой
квадратичной функцией с коэффициентом $3R^2+E^2>0$.
При $DR\ne0$ она имеет единственный абсолютный максимум на
$[-1,1]$, что противоречит наличию двух различных точек касания.
Значит, $E=0$, и два максимума существуют при
\begin {equation}
 F^2>3R^2,\qquad
 c_*:=\cos\psi_*=\frac{{2DR}}{{F^2-3R^2}},
 \qquad |c_*|<1.
 \label {eq:ndc-primer-angles}
\end {equation}
Эти формулы получаются дифференцированием
$f(c)=D^2+R^2+F^2-(F^2-3R^2)c^2+4DRc$.
Условие нормировки имеет вид
\begin {equation}
 D^2+R^2+F^2+\frac{{4D^2R^2}}{{F^2-3R^2}}=1.
 \label {eq:ndc-primer-normalization}
\end {equation}
В двух точках касания трансверсальные компоненты базис-вектора
одинаковы, а радиальные и нормальные компоненты противоположны.
При $|c_*|=1$ точки сливаются; это граничный случай с возможным
единственным ненулевым импульсом.

\paragraph{Окружность: сингулярный случай (SC).}
Если центральный эллипс имеет постоянный единичный радиус,
разделение на отдельные точки максимума исчезает.
Из $D=0$ и $Q$ равной единичной матрице получаем
\begin {equation}
 \begin {gathered}
 E=0,\qquad R=\frac{1}{2},\qquad F=\pm\frac{{\sqrt3}}{{2}},\\
 \boldsymbol p(\psi)=
 \begin {pmatrix}
   \tfrac{1}{2}\sin\psi\\ \cos\psi\\
   \pm\tfrac{{\sqrt3}}{{2}}\sin\psi
 \end {pmatrix},
 \qquad \lVert\boldsymbol p(\psi)\rVert\equiv1.
 \end {gathered}
 \label {eq:singular-primer-circle}
\end {equation}
Все фазы имеют одинаковую предельную эффективность, но направления
импульсов по-прежнему заданы функцией $\boldsymbol p(\psi)$.
Свобода выбора фаз не снимает граничных условий и требования
неотрицательности величин ${\Delta\widehat V_j}$.
Различные допустимые распределения импульсов, удовлетворяющие
\eqref{eq:certificate}, имеют одинаковую оптимальную стоимость.
В частности, постоянным здесь является модуль, а не сам вектор.

\paragraph{Вырождение в отрезок и точку.}
При $R=0$ угол $\alpha$ не определён; делить на $R$ нельзя.
В исходных координатах
\begin {equation}
 \boldsymbol p(\varphi)=
      (0,D,l_s\sin\varphi+l_c\cos\varphi)^{\mathsf T}.
 \label {eq:line-primer-general}
\end {equation}
Если $l_s^2+l_c^2>0$, годограф --- отрезок, и
$\max\lVert\boldsymbol p\rVert^2=D^2+l_s^2+l_c^2$.
Два конца отрезка соответствуют противоположным фазам;
трансверсальные компоненты одинаковы, нормальные противоположны,
радиальные равны нулю.
При $l_s=l_c=0$ годограф вырождается в точку, а нормировка даёт
$D=\pm1$. Этот предел особенно существенен для компланарных
переходов с преобладающим изменением большой полуоси.

Таким образом, геометрия базис-вектора определяет допустимые типы
расположения импульсов. Она ещё не определяет, какой тип отвечает
заданным ${\Delta a},{\Delta e_x},{\Delta e_y},{\Delta i}$: это устанавливается совместным решением
граничных уравнений, условий касания и неравенств на величины
импульсов. Именно эта задача решается в следующем разделе.

% ==================== planar ====================
\section{Алгоритмы перехода}
\label{sec:algorithms}

Исходными данными служат пять изменений орбитальных параметров из
\eqref{eq:boundary-map}. После выбора оси вдоль линии взаимных узлов
модуль изменения вектора эксцентриситета обозначим через $\Delta e$:
\begin {equation}
 \Delta e=\sqrt{{(\Delta e_x)}^2+{(\Delta e_y)}^2}.
 \label {eq:short-data}
\end {equation}
Величина $\Delta e$ равна модулю разности векторов эксцентриситета;
в общем случае она не равна разности их модулей $e_f-e_i$.
Величина $\Delta a$ безразмерна.
Искомые компоненты ${\Delta\widehat V_{rj}},{\Delta\widehat V_{tj}},{\Delta\widehat V_{zj}}$ также безразмерны;
окончательный переход к размерным импульсам выполняется умножением на $V_0$.
Во всех формулах ниже ${\Delta\widehat V_j}=\lVert\Delta\widehat{\boldsymbol V}_j\rVert\geqslant0$.
Знак отдельной компоненты импульса не является знаком его модуля.

\subsection{Компланарный случай}
\label{sec:planar}

\subsubsection{Нижняя оценка характеристической скорости}

При ${\Delta i}=0$ нормальные компоненты не нужны. Действительно, их удаление
сохраняет три плоских граничных условия и не увеличивает стоимость.
Введём комплексную запись разности эксцентриситетов:
\begin {equation}
 {\Delta e_x}+\mathrm i {\Delta e_y}
 =\sum_j(2{\Delta\widehat V_{tj}}-\mathrm i {\Delta\widehat V_{rj}})\exp(\mathrm i\varphi_j),
 \qquad {\Delta a}=2\sum_j{\Delta\widehat V_{tj}}.
 \label {eq:planar-complex}
\end {equation}
Мнимая единица $\mathrm i$ в этой формуле не обозначает наклонение орбиты.
Модуль комплексного множителя одного импульса равен
$\sqrt{4{(\Delta\widehat V_{tj})}^2+{(\Delta\widehat V_{rj})}^2}$ и не превосходит $2{\Delta\widehat V_j}$.
По неравенству треугольника
\[
 |{\Delta a}|\leqslant2\sum_j|{\Delta\widehat V_{tj}}|\leqslant2J,
 \qquad
 \Delta e\leqslant\sum_j\sqrt{4{(\Delta\widehat V_{tj})}^2+{(\Delta\widehat V_{rj})}^2}\leqslant2J.
\]
Следовательно, для любого числа импульсов
\begin {equation}
 J\geqslant\frac{1}{2}\max\{|{\Delta a}|,\Delta e\}.
 \label {eq:planar-bound}
\end {equation}
Ниже построим управление, достигающее этой оценки. Поэтому оценка
сразу станет формулой глобального минимума линейной задачи.

\subsubsection{Два противоположных места приложения}

Пусть $\Delta e>0$ и
\begin {equation}
 \varphi_e=\operatorname {atan2}({\Delta e_y},{\Delta e_x}),\qquad
 \varphi_1=\varphi_e,\qquad\varphi_2=\varphi_e+\pi.
 \label {eq:planar-locations}
\end {equation}
Функция $\operatorname {atan2}$ восстанавливает угол по двум компонентам
с учётом квадранта. Замена её на $\arctan({\Delta e_y}/{\Delta e_x})$ без проверки квадранта
может переставить разгонный и тормозной импульсы.
Выберем оба импульса трансверсальными. Граничные условия принимают вид
\[
 2({\Delta\widehat V_{t1}}+{\Delta\widehat V_{t2}})={\Delta a},\qquad
 2({\Delta\widehat V_{t1}}-{\Delta\widehat V_{t2}})=\Delta e.
\]
Отсюда
\begin {equation}
 \boxed{\quad {\Delta\widehat V_{t1}}=\frac{{{\Delta a}+\Delta e}}{{4}},\qquad
 {\Delta\widehat V_{t2}}=\frac{{{\Delta a}-\Delta e}}{{4}},\qquad {\Delta\widehat V_{r1}}={\Delta\widehat V_{r2}}=0.\quad}
 \label {eq:planar-impulses}
\end {equation}
Стоимость равна
\begin {equation}
 \boxed{\quad J_{\rm pl}=\frac{{|{\Delta a}+\Delta e|+|{\Delta a}-\Delta e|}}{{4}}
 =\frac{1}{2}\max\{|{\Delta a}|,\Delta e\}.\quad}
 \label {eq:planar-cost}
\end {equation}
Равенство с нижней оценкой доказывает оптимальность, включая возможность
любого другого числа импульсов. Положительная ${\Delta\widehat V_t}$ означает разгон по
трансверсальному направлению, отрицательная --- торможение.

Приведённая пара углов относится к оси \emph{разности} векторов
эксцентриситета. Она не обязана совпадать с линией апсид ни начальной,
ни конечной орбиты. Это существенно при различных ориентациях их перицентров.

\subsubsection{Пересекающиеся и непересекающиеся орбиты}

В первом порядке радиус орбиты с параметрами $(a,e_x,e_y)$ равен
\[
 r(\varphi)=a\,[1-e_x\cos\varphi-e_y\sin\varphi]+O(\varepsilon^2r_0).
\]
Поэтому разность радиусов конечной и начальной орбит выражается как
\begin {equation}
 \frac{{r_f(\varphi)-r_i(\varphi)}}{{r_0}}
 ={\Delta a}-{\Delta e_x}\cos\varphi-{\Delta e_y}\sin\varphi
 ={\Delta a}-\Delta e\cos(\varphi-\varphi_e).
 \label {eq:planar-intersection}
\end {equation}
В этой формуле и далее классификация пересечения относится именно
к геометрии первого порядка.

Если $\Delta e>|{\Delta a}|$, уравнение равенства радиусов имеет два решения.
В \eqref{eq:planar-impulses} первая компонента положительна, вторая
отрицательна; стоимость определяется изменением эксцентриситета.
Равенство в цепочке оценок для $\Delta e$ требует ${\Delta\widehat V_{rj}}=0$ и сонаправленности
всех ненулевых вкладов в разность векторов эксцентриситета.
Следовательно, положительные трансверсальные импульсы могут располагаться
только на луче $\varphi_e$, а отрицательные --- на противоположном луче.
Импульсы, приложенные в одном месте на разных оборотах, можно объединить.
Таким образом, два места \eqref{eq:planar-locations} не являются произвольным
выбором: они следуют из условий достижения минимальной стоимости.

Если $|{\Delta a}|>\Delta e$, орбиты не пересекаются. Теперь стоимость определяется
изменением большой полуоси. Равенство $|{\Delta a}|=2J$ требует, чтобы все импульсы
были трансверсальными и имели один знак $\operatorname {sgn}{\Delta a}$.
Их места не определяются однозначно. При двух импульсах положим
\[
 w_j=\frac{{2{\Delta\widehat V_{tj}}}}{{{\Delta a}}}\geqslant0,\qquad w_1+w_2=1,
 \qquad \boldsymbol m=\frac{1}{{\Delta a}}\begin {pmatrix}{\Delta e_x}\\{\Delta e_y}\end {pmatrix}.
\]
Тогда оставшиеся условия равносильны
\begin {equation}
 \boldsymbol m=w_1\boldsymbol u_1+w_2\boldsymbol u_2,
 \qquad \boldsymbol u_j=(\cos\varphi_j,\sin\varphi_j)^\mathsf T.
 \label {eq:planar-chord}
\end {equation}
Поскольку $\lVert\boldsymbol m\rVert=\Delta e/|{\Delta a}|<1$, точка
$\boldsymbol m$ лежит внутри единичного круга. Любая хорда, проходящая
через неё, даёт допустимую пару импульсов той же стоимости.

\begin{figure}[htbp]
\centering
\begin{tikzpicture}[scale=1.05,>=Latex]
 \draw[->,gray] (-2.15,0)--(2.25,0) node[right] {$u_x=\cos\varphi$};
 \draw[->,gray] (0,-1.95)--(0,2.15) node[above] {$u_y=\sin\varphi$};
 \draw[accent,thick] (0,0) circle(1.7);
 \coordinate (U1) at (1.597,.581);
 \coordinate (U2) at (-1.303,-1.092);
 \coordinate (M) at (.437,-.088);
 \draw[thick] (U2)--(U1);
 \fill[accent] (U1) circle(1.7pt) node[above right] {$\boldsymbol u_1$};
 \fill[accent] (U2) circle(1.7pt) node[below left] {$\boldsymbol u_2$};
 \fill (M) circle(1.7pt) node[below right] {$\boldsymbol m$};
 \draw[->,accent] (0,0)--(U1);
 \draw[->,accent] (0,0)--(U2);
\end{tikzpicture}
\caption{Два трансверсальных импульса одного знака. Заданный вектор
$\boldsymbol m$ представляется выпуклой комбинацией концов хорды.
Вес определяет долю изменения большой полуоси.}
\label{fig:planar-chord}
\end{figure}

Для явного расчёта выберем произвольное единичное направление
$\boldsymbol v$ хорды и найдём пересечения прямой
$\boldsymbol m+t\boldsymbol v$ с окружностью:
\begin {equation}
 t_\pm=-\boldsymbol m\cdot\boldsymbol v
 \pm\sqrt{(\boldsymbol m\cdot\boldsymbol v)^2+1-\lVert\boldsymbol m\rVert^2}.
 \label {eq:chord-roots}
\end {equation}
Здесь $t_-<0<t_+$. Полагая
\[
 \boldsymbol u_\pm=\boldsymbol m+t_\pm\boldsymbol v,\qquad
 w_+=\frac{{-t_-}}{{t_+-t_-}},\qquad
 w_-=\frac{{t_+}}{{t_+-t_-}},
\]
получаем ${\Delta\widehat V_{t\pm}}={\Delta a}w_\pm/2$ и
$\varphi_\pm=\operatorname {atan2}(u_{\pm,y},u_{\pm,x})$.
Оба веса положительны. Формулы \eqref{eq:planar-impulses}
дают частную, наиболее простую для расчёта пару.

\subsubsection{Границы и нулевые данные}

При $\Delta e=|{\Delta a}|>0$ один из двух импульсов
\eqref{eq:planar-impulses} обращается в нуль.
Остаётся один трансверсальный импульс в точке касания орбит первого порядка.
При ${\Delta a}>0$ это $\varphi_e$; при ${\Delta a}<0$ это $\varphi_e+\pi$.
Не следует отождествлять это утверждение с существованием точного касания
двух заданных кеплеровых эллипсов: малые члены второго порядка могут
изменить их взаимное положение.

Если $\Delta e=0$, угол $\varphi_e$ не определён и вычислять его не нужно.
При ${\Delta a}\ne0$ выбирают любую пару противоположных углов и
${\Delta\widehat V_{t1}}={\Delta\widehat V_{t2}}={\Delta a}/4$. Это линейный аналог перехода между близкими
круговыми орбитами. При ${\Delta a}=\Delta e=0$ манёвры не требуются.
Случай ${\Delta a}=0$, $\Delta e>0$ даёт два противоположных трансверсальных
импульса величин $\Delta e/4$ каждого знака; суммарное изменение
большой полуоси равно нулю, но промежуточная орбита отличается от исходной.

\subsubsection{Расчётный пример}

Пусть $r_0=7000$ км, $\mu=398600{,}4418$ км$^3$/с$^2$,
${\Delta a}=0{,}001$, ${\Delta e_x}=0{,}002$, ${\Delta e_y}=0{,}001$.
Тогда $V_0=7{,}54605329$ км/с и
\[
 \Delta e=0{,}00223606798>|{\Delta a}|,\qquad
 \varphi_1=26{,}565051^\circ,\qquad\varphi_2=206{,}565051^\circ.
\]
Из \eqref{eq:planar-impulses}
\[
 {\Delta\widehat V_{t1}}=0{,}000809016994,\qquad {\Delta\widehat V_{t2}}=-0{,}000309016994,
 \qquad J=0{,}001118033989.
\]
Размерные компоненты равны приблизительно $6{,}1049$ и $-2{,}3319$ м/с,
суммарная характеристическая скорость --- $8{,}4367$ м/с.
Проверка изменения большой полуоси даёт
$2({\Delta\widehat V_{t1}}+{\Delta\widehat V_{t2}})=0{,}001$; проверка эксцентриситета ---
$2({\Delta\widehat V_{t1}}-{\Delta\widehat V_{t2}})(\cos\varphi_e,\sin\varphi_e)=(0{,}002,0{,}001)$.
Разность знаков и разность модулей имеют разный смысл:
первая создаёт нужное изменение вектора эксцентриситета,
вторая оставляет ненулевое суммарное изменение большой полуоси.

Полученные плоские решения совпадают с классификацией
\cite[гл.~2]{baranov-ru,baranov-en}. Их оптимальность здесь доказана
непосредственно оценкой \eqref{eq:planar-bound}, поэтому не зависит
от выбранной параметризации базис-вектора.

% ==================== spatial ====================
\subsection{Некомпланарный случай}
\label{sec:spatial}

\subsubsection{Подготовка исходных данных}

Пусть $\boldsymbol k_i,\boldsymbol k_f$ --- единичные векторы
орбитальных моментов начальной и конечной орбит.
При ненулевом малом угле между плоскостями построим базис
\begin {equation}
 \boldsymbol e_X=\frac{{\boldsymbol k_i\times\boldsymbol k_f}}{{\lVert\boldsymbol k_i\times\boldsymbol k_f\rVert}},\qquad
 \boldsymbol e_Z=\boldsymbol k_i,\qquad
 \boldsymbol e_Y=\boldsymbol e_Z\times\boldsymbol e_X.
 \label {eq:mutual-nodes-basis}
\end {equation}
Ось $X$ направлена к восходящему узлу конечной плоскости относительно
начальной. Угол $\varphi$ отсчитывается от неё в направлении движения
по начальной орбите. Если $\boldsymbol e_i,\boldsymbol e_f$ ---
пространственные векторы эксцентриситета, то
\begin {equation}
 {\Delta e_x}=(\boldsymbol e_f-\boldsymbol e_i)\cdot\boldsymbol e_X,\qquad
 {\Delta e_y}=(\boldsymbol e_f-\boldsymbol e_i)\cdot\boldsymbol e_Y,\qquad
 {\Delta i}=\operatorname {atan2}\bigl(\lVert\boldsymbol k_i\times\boldsymbol k_f\rVert,
 \boldsymbol k_i\cdot\boldsymbol k_f\bigr).
 \label {eq:spatial-input}
\end {equation}
Проекция векторов эксцентриситета на общую плоскость согласована
с первым порядком: возникающие из-за взаимного наклона поправки имеют
высший порядок. Для плоской задачи ${\Delta i}=0$ узловой базис не определён,
но он и не требуется: применяется раздел~\ref{sec:planar} с любой
фиксированной осью в общей плоскости.

Если исходные данные заданы кеплеровыми элементами, взаимный угол
следует вычислять по
\[
 \cos {\Delta i}=\cos i_i\cos i_f+\sin i_i\sin i_f\cos(\Omega_f-\Omega_i),
\]
либо по более устойчивой вблизи нуля формуле \eqref{eq:spatial-input}.
В общем случае ${\Delta i}\ne|i_f-i_i|$. В обозначениях Баранова встречается
также величина $2\sin({\Delta i}/2)$; её отличие от ${\Delta i}$ имеет третий порядок
и не влияет на рассматриваемую линейную теорию.

\subsubsection{Области применимости}

Положим
\begin {equation}
 B_* ={(\Delta e_x)}^2+{(\Delta e_y)}^2+\frac{{2}}{{\sqrt3}}{\Delta i}{|\Delta e_y|}-{(\Delta i)}^2.
 \label {eq:classification-B}
\end {equation}
Буква $B_*$ обозначает скаляр, в отличие от матрицы $B(\varphi)$.
Классификация имеет вид
\begin {align}
 \mathrm {NC}:&\quad {(\Delta a)}^2\leqslant {(\Delta e_x)}^2,\qquad {(\Delta i)}^2\geqslant3{(\Delta e_y)}^2,
 \label {eq:NC-region}\\
 \mathrm {SC}:&\quad {(\Delta a)}^2\leqslant B_*,\qquad {(\Delta i)}^2\leqslant3{(\Delta e_y)}^2,
 \label {eq:SC-region-summary}\\
 \mathrm {NDC}:&\quad {(\Delta a)}^2\geqslant\max\{{(\Delta e_x)}^2,B_*\}.
 \label {eq:NDC-region}
\end {align}
Здесь NC --- узловое решение, NDC --- невырожденное пространственное
решение, SC --- особое, или сингулярное, решение.
Области покрывают все данные. Действительно,
$B_*-{(\Delta e_x)}^2=({|\Delta e_y|}-{\Delta i}/\sqrt3)({|\Delta e_y|}+\sqrt3{\Delta i})$:
при ${\Delta i}\geqslant\sqrt3{|\Delta e_y|}$ имеем $B_*\leqslant {(\Delta e_x)}^2$, а при обратном
неравенстве --- $B_*\geqslant {(\Delta e_x)}^2$. На общих границах стоимости
соседних решений совпадают; один набор данных может допускать несколько
равноценных представлений управления.

Ниже области \eqref{eq:NC-region}--\eqref{eq:NDC-region}
будут получены из условий существования допустимого базис-вектора
и неотрицательности модулей импульсов. Их нельзя использовать независимо
от этих условий, например выбирая просто наименьшее число из трёх
формул стоимости вне соответствующих областей.

\subsubsection{Узловое решение NC}
\label{sec:NC}

В этом случае импульсы располагаются при $\varphi_1=0$ и
$\varphi_2=\pi$. Их единичные направления в локальных осях
противоположны. Обозначим первое направление через
$\boldsymbol d=(d_r,d_t,d_z)^\mathsf T$ и положим
$\Delta\widehat{\boldsymbol V}_1={\Delta\widehat V_1}\boldsymbol d$,
$\Delta\widehat{\boldsymbol V}_2=-{\Delta\widehat V_2}\boldsymbol d$.
Подстановка в граничные условия даёт
\[
 \begin{aligned}
 \Delta a&=2d_t(\Delta\widehat V_1-\Delta\widehat V_2),&
 \Delta e_x&=2d_t(\Delta\widehat V_1+\Delta\widehat V_2),\\
 \Delta e_y&=-d_r(\Delta\widehat V_1+\Delta\widehat V_2),&
 \Delta i&=d_z(\Delta\widehat V_1+\Delta\widehat V_2).
 \end{aligned}
\]
Поскольку ${\Delta\widehat V_1}+{\Delta\widehat V_2}=J$ и $\lVert\boldsymbol d\rVert=1$,
\begin {equation}
 \boxed{J_{\rm NC}=\sqrt{{(\Delta i)}^2+\frac{{{(\Delta e_x)}^2}}{{4}}+{(\Delta e_y)}^2},\qquad
 \boldsymbol d=\frac{1}{{J_{\rm NC}}}\begin {pmatrix}-{\Delta e_y}\\{\Delta e_x}/2\\{\Delta i}\end {pmatrix}.}
 \label {eq:NC-cost}
\end {equation}
При ${\Delta e_x}\ne0$ модули равны
\begin {equation}
 {\Delta\widehat V_1}=\frac{{J_{\rm NC}}}{2}\left(1+\frac{{\Delta a}}{{\Delta e_x}}\right),\qquad
 {\Delta\widehat V_2}=\frac{{J_{\rm NC}}}{2}\left(1-\frac{{\Delta a}}{{\Delta e_x}}\right).
 \label {eq:NC-magnitudes}
\end {equation}
Оба модуля неотрицательны тогда и только тогда, когда ${(\Delta a)}^2\leqslant {(\Delta e_x)}^2$.
Знак ${\Delta e_x}$ сохраняется в формулах: заменять ${\Delta a}/{\Delta e_x}$ на $|{\Delta a}/{\Delta e_x}|$ нельзя.

Остаётся доказать оптимальность среди произвольных, в том числе
неузловых, схем. Для этого выберем множители
\begin {equation}
 \boldsymbol l=\frac{1}{{J_{\rm NC}}}
 \begin {pmatrix}0\\{\Delta e_x}/4\\{\Delta e_y}\\-3{\Delta e_x}{\Delta e_y}/(4{\Delta i})\\{\Delta i}\end {pmatrix}.
 \label {eq:NC-dual}
\end {equation}
Непосредственная подстановка в \eqref{eq:primer-components} даёт
$\boldsymbol p(0)=\boldsymbol d$,
$\boldsymbol p(\pi)=-\boldsymbol d$ и
\begin {equation}
 \lVert\boldsymbol p(\varphi)\rVert^2
 =1-\frac{{({(\Delta i)}^2-3{(\Delta e_y)}^2)(1+3{(\Delta e_x)}^2/(16{(\Delta i)}^2))}}{{J_{\rm NC}^2}}
 \sin^2\varphi.
 \label {eq:NC-primer-norm}
\end {equation}
Условие ${(\Delta i)}^2\geqslant3{(\Delta e_y)}^2$ обеспечивает
$\lVert\boldsymbol p\rVert\leqslant1$ на всём обороте.
Равенство достигается в местах импульсов; следовательно, выполняется
критерий достаточности \eqref{eq:certificate}.
Это выводит вторую границу области NC и доказывает глобальный минимум
линейной задачи.

Если ${\Delta e_x}=0$, условие области требует ${\Delta a}=0$.
Формулы \eqref{eq:NC-magnitudes} применять нельзя, однако сама задача
остаётся определённой. Достаточен один импульс в узле $\varphi=0$:
$\Delta\widehat{\boldsymbol V}=(-{\Delta e_y},0,{\Delta i})^\mathsf T$.
Можно также произвольно разделить его модуль между противоположными
узлами с направлениями $\boldsymbol d$ и $-\boldsymbol d$.
На границе ${(\Delta i)}^2=3{(\Delta e_y)}^2$ норма \eqref{eq:NC-primer-norm}
постоянна, и NC соприкасается с особым решением.

\paragraph{Пример.}
Пусть ${\Delta a}=0{,}001$, ${\Delta e_x}=0{,}002$, ${\Delta e_y}=0{,}0005$, ${\Delta i}=0{,}002$.
Оба условия NC выполнены. Получаем
\[
 J=0{,}002291287847,\qquad {\Delta\widehat V_1}=0{,}001718465886,
 \qquad {\Delta\widehat V_2}=0{,}000572821962,
\]
\begin {equation}
 \Delta\widehat{\boldsymbol V}_1=\begin {pmatrix}-0{,}000375\\0{,}000750\\0{,}001500\end {pmatrix},
 \qquad
 \Delta\widehat{\boldsymbol V}_2=\begin {pmatrix}0{,}000125\\-0{,}000250\\-0{,}000500\end {pmatrix}.
 \label {eq:NC-example}
\end {equation}
Например, ${\Delta e_x}=2({\Delta\widehat V_{t1}}-{\Delta\widehat V_{t2}})=0{,}002$,
${\Delta e_y}=-{\Delta\widehat V_{r1}}+{\Delta\widehat V_{r2}}=0{,}0005$,
${\Delta i}={\Delta\widehat V_{z1}}-{\Delta\widehat V_{z2}}=0{,}002$.
При $r_0=7000$ км размерная стоимость равна приблизительно $17{,}2902$ м/с.

\subsubsection{Невырожденное решение NDC: вывод стоимости}
\label{sec:NDC}

Параметрическое представление этого семейства используется
в~\cite[с.~48--53]{edelbaum} и в линейной части~\cite{gobetz}.
Сначала рассмотрим внутренность области NDC при ${\Delta i}>0$, ${|\Delta e_y|}>0$.
Для двух симметричных максимумов базис-вектор имеет вид
\begin {equation}
 \boldsymbol p(\varphi)=
 \begin {pmatrix}R\sin\psi\\D+2R\cos\psi\\F\sin\psi\end {pmatrix},
 \qquad \psi=\varphi-\alpha,
 \qquad R>0,\quad F^2>3R^2.
 \label {eq:NDC-primer}
\end {equation}
Буква $D$ здесь обозначает постоянную трансверсальную составляющую
базис-вектора, а не характеристическую скорость.
Если $c=\cos\psi$, то
\[
 p^2=R^2+D^2+F^2+4RDc-(F^2-3R^2)c^2.
\]
Это вогнутая квадратичная функция $c$.
Её максимум достигается при
\begin {equation}
 c_0=\frac{{2RD}}{{F^2-3R^2}},\qquad
 (R^2+F^2)\left(1+\frac{{D^2}}{{F^2-3R^2}}\right)=1.
 \label {eq:NDC-normalization}
\end {equation}
Последнее равенство задаёт единичное значение максимума; дополнительно
требуется $|c_0|\leqslant1$.

Построение оптимального базис-вектора удобно выполнять как максимизацию
нижней оценки $\boldsymbol l\cdot\boldsymbol b$.
Соответствие \eqref{eq:NDC-primer} компонентам множителей даёт
\begin {equation}
 \boldsymbol l\cdot\boldsymbol b
 =\frac{{{\Delta a}D}}{{2}}+R{\Delta e_x}\cos\alpha+(R{\Delta e_y}-{\Delta i}F)\sin\alpha.
 \label {eq:NDC-dual-objective}
\end {equation}
При фиксированных $R,D,F$ максимум по $\alpha$ равен
${\Delta a}D/2+\sqrt{R^2{(\Delta e_x)}^2+(R{\Delta e_y}-{\Delta i}F)^2}$.
Для $R>0$ выбираем знак $F$, противоположный знаку ${\Delta e_y}$;
тогда $|R{\Delta e_y}-{\Delta i}F|=R|{\Delta e_y}|+{\Delta i}|F|$.
Введём
\[
 \eta=\frac{{|F|}}{{R}}>\sqrt3,\qquad
 k=\sqrt{R^2+F^2},\qquad0\leqslant k\leqslant1.
\]
Из \eqref{eq:NDC-normalization}
\[
 |D|=\sqrt{1-k^2}\sqrt{\frac{{\eta^2-3}}{{1+\eta^2}}},\qquad
 R=\frac{{k}}{{\sqrt{1+\eta^2}}}.
\]
Знак $D$ выбираем по знаку ${\Delta a}$.
Поэтому выражение \eqref{eq:NDC-dual-objective}, максимизированное
по ориентации, принимает вид
\[
 a_\eta\sqrt{1-k^2}+b_\eta k,\qquad
 a_\eta=\frac{{|{\Delta a}|}}{{2}}\sqrt{\frac{{\eta^2-3}}{{1+\eta^2}}},\qquad
 b_\eta=\frac{{\sqrt{{(\Delta e_x)}^2+({|\Delta e_y|}+{\Delta i}\eta)^2}}}{{\sqrt{1+\eta^2}}}.
\]
Неравенство Коши даёт максимум $\sqrt{a_\eta^2+b_\eta^2}$,
достигаемый при $k=b_\eta/\sqrt{a_\eta^2+b_\eta^2}$.
Квадрат этой величины равен
\begin {equation}
 \frac{{({(\Delta i)}^2+{(\Delta a)}^2/4)\eta^2+2{\Delta i}{|\Delta e_y|}\eta+({(\Delta e_x)}^2+{(\Delta e_y)}^2-3{(\Delta a)}^2/4)}}{{1+\eta^2}}.
 \label {eq:NDC-rayleigh}
\end {equation}
Это отношение $\boldsymbol v^\mathsf T M_N\boldsymbol v/\lVert\boldsymbol v\rVert^2$
при $\boldsymbol v=(\eta,1)^\mathsf T$ и
\begin {equation}
 M_N=\begin {pmatrix}
 {(\Delta i)}^2+{(\Delta a)}^2/4&{\Delta i}{|\Delta e_y|}\\{\Delta i}{|\Delta e_y|}&{(\Delta e_x)}^2+{(\Delta e_y)}^2-3{(\Delta a)}^2/4
 \end {pmatrix}.
 \label {eq:NDC-matrix}
\end {equation}
Максимум отношения для вещественной симметричной матрицы равен её
наибольшему собственному значению. Это следует, например, из разложения
по ортонормированным собственным векторам: отношение становится
взвешенным средним собственных значений с неотрицательными весами.

Для матрицы $2\times2$ собственные значения находятся из квадратного
уравнения. Обозначив
\begin {equation}
 K={(\Delta i)}^2-{(\Delta e_x)}^2-{(\Delta e_y)}^2+{(\Delta a)}^2,\qquad H=\sqrt{K^2+4{(\Delta i)}^2{(\Delta e_y)}^2},
 \label {eq:NDC-KH}
\end {equation}
получаем
\begin {equation}
 \boxed{\quad J_{\rm NDC}=
 \sqrt{\frac{{{(\Delta i)}^2+{(\Delta e_x)}^2+{(\Delta e_y)}^2-{(\Delta a)}^2/2+H}}{{2}}}.\quad}
 \label {eq:NDC-cost}
\end {equation}
Пока найдена кандидатная нижняя оценка в выбранном семействе.
Для доказательства, что это стоимость допустимого оптимального перехода,
необходимо восстановить импульсы, проверить ограничение $|c_0|\leqslant1$
и неотрицательность их модулей. Эти проверки выполняются далее.

\subsubsection{NDC: углы, направления и модули импульсов}

Собственный вектор, отвечающий большему собственному значению
\eqref{eq:NDC-matrix}, имеет положительные компоненты при ${\Delta i}{|\Delta e_y|}>0$.
Удобно вычислять его отношение
\begin {equation}
 \eta=\frac{{K+H}}{{2{\Delta i}{|\Delta e_y|}}},\qquad
 L=\sqrt{{(\Delta e_x)}^2+({|\Delta e_y|}+{\Delta i}\eta)^2},\qquad \sigma=\operatorname {sgn}{\Delta e_y}.
 \label {eq:NDC-eta}
\end {equation}
Символ $L$ здесь вспомогательный; он не обозначает функцию Лагранжа.
Восстановление параметров базис-вектора даёт
\begin {align}
 \alpha&=\operatorname {atan2}\bigl(\sigma({|\Delta e_y|}+{\Delta i}\eta),{\Delta e_x}\bigr),
 \label {eq:NDC-alpha}\\
 R&=\frac{{L}}{{(1+\eta^2)J}},\qquad
 D=\frac{{{\Delta a}(\eta^2-3)}}{{2(1+\eta^2)J}},\qquad F=-\sigma\eta R.
 \label {eq:NDC-RDF}
\end {align}
Здесь и до конца этого подраздела $J=J_{\rm NDC}$.
Из \eqref{eq:NDC-normalization} следует $c_0={\Delta a}/L$, поэтому
\begin {equation}
 \theta=\arccos({\Delta a}/L),\qquad
 \varphi_+=\alpha+\theta,\qquad\varphi_-=\alpha-\theta.
 \label {eq:NDC-angles}
\end {equation}
Индексы $+$ и $-$ обозначают две геометрические ветви и не задают
хронологический порядок манёвров.
Единичные направления равны
\begin {equation}
 \boldsymbol d_\pm=
 \begin {pmatrix}\pm R\sin\theta\\D+2R\cos\theta\\\pm F\sin\theta\end {pmatrix}.
 \label {eq:NDC-directions}
\end {equation}
Радиальные и нормальные компоненты противоположны, трансверсальные
совпадают. Это утверждение касается направлений, а не размерных
компонент: модули импульсов, вообще говоря, различны.

Положим ${\Delta\widehat V_\pm}=J(1\pm\delta)/2$. Для вывода $\delta$ запишем
нормальные граничные условия в комплексной форме:
\begin {equation}
 {\Delta i}=\sum_j{\Delta\widehat V_{zj}}\exp(\mathrm i\varphi_j)
 =JF\sin\theta\,\exp(\mathrm i\alpha)
 (\delta\cos\theta+\mathrm i\sin\theta).
 \label {eq:NDC-plane-balance}
\end {equation}
Мнимая часть равна нулю, откуда
$\delta\cos\theta\sin\alpha+\sin\theta\cos\alpha=0$.
Подстановка \eqref{eq:NDC-alpha} и \eqref{eq:NDC-angles} даёт
\begin {equation}
 \boxed{\quad
 \delta=-\frac{{\sigma {\Delta e_x}\sqrt{L^2-{(\Delta a)}^2}}}{{{\Delta a}({|\Delta e_y|}+{\Delta i}\eta)}},\qquad
 {\Delta\widehat V_\pm}=\frac{J}{2}(1\pm\delta),\qquad
 \Delta\widehat{\boldsymbol V}_\pm={\Delta\widehat V_\pm}\boldsymbol d_\pm.\quad}
 \label {eq:NDC-magnitudes}
\end {equation}
Во внутренности области NDC ${\Delta a}\ne0$, поэтому знаменатель не обращается
в нуль. Граничные и вырожденные данные рассматриваются отдельно.

Проверим область применимости. Из уравнения собственного вектора
матрицы \eqref{eq:NDC-matrix}
\begin {equation}
 K={\Delta i}{|\Delta e_y|}(\eta-\eta^{-1}).
 \label {eq:NDC-eigen-identity}
\end {equation}
Правая часть строго возрастает при $\eta>0$.
Поэтому $\eta>\sqrt3$ равносильно ${(\Delta a)}^2>B_*$.
Ещё одно следствие \eqref{eq:NDC-eigen-identity}:
\begin {equation}
 L^2-{(\Delta a)}^2=\frac{{{\Delta i}(1+\eta^2)({|\Delta e_y|}+{\Delta i}\eta)}}{{\eta}}>0.
 \label {eq:NDC-cos-feasible}
\end {equation}
Таким образом, угол \eqref{eq:NDC-angles} существует и $|c_0|<1$.
Наконец,
\[
 \delta^2\leqslant1
 \quad\Longleftrightarrow\quad
 {(\Delta e_x)}^2(L^2-{(\Delta a)}^2)\leqslant {(\Delta a)}^2({|\Delta e_y|}+{\Delta i}\eta)^2
 \quad\Longleftrightarrow\quad {(\Delta e_x)}^2\leqslant {(\Delta a)}^2.
\]
Это второе условие области NDC. При строгом неравенстве оба модуля
положительны; при равенстве один из них может исчезнуть.

Для полноты проверим остальные граничные условия.
В точках максимумов $D+2R\cos\theta={\Delta a}/(2J)$, поэтому
$2({\Delta\widehat V_{t+}}+{\Delta\widehat V_{t-}})={\Delta a}$.
Комплексная разность эксцентриситетов вычисляется как
\begin {equation}
 {\Delta e_x}+\mathrm i {\Delta e_y}=
 J\exp(\mathrm i\alpha)
 \left[2d_t\cos\theta+R\sin^2\theta
 +\mathrm i\delta(2d_t\sin\theta-R\sin\theta\cos\theta)\right],
 \label {eq:NDC-ecc-balance}
\end {equation}
где $d_t={\Delta a}/(2J)$. После подстановки
\eqref{eq:NDC-eta}--\eqref{eq:NDC-magnitudes} действительная
и мнимая части равны соответственно ${\Delta e_x}$ и ${\Delta e_y}$; сокращения используют
тождество \eqref{eq:NDC-eigen-identity}.
Действительная часть \eqref{eq:NDC-plane-balance} равна ${\Delta i}$
в силу \eqref{eq:NDC-cos-feasible}.
Тем самым проверены все пять условий.

В силу \eqref{eq:NDC-normalization} базис-вектор имеет норму не более
единицы на всём обороте, а в двух выбранных точках его норма равна единице.
Модули неотрицательны. Критерий достаточности доказывает, что
\eqref{eq:NDC-cost} является минимумом полной линейной задачи,
а не только минимумом в заранее выбранном двухимпульсном классе.

\paragraph{Замечание о вычислениях.}
При очень малом ${|\Delta e_y|}$ отношение $\eta$ может быть велико.
В таком случае численно устойчивее вычислить нормированный собственный
вектор матрицы \eqref{eq:NDC-matrix}, чем образовывать большие степени
$\eta$. Значение ${|\Delta e_y|}=0$ следует обрабатывать отдельными формулами ниже.
Округление аргумента $\arccos$ до $[-1,1]$ допустимо лишь для устранения
ошибки машинного округления после проверки области применимости.

\paragraph{Пример.}
При ${\Delta a}=0{,}004$, ${\Delta e_x}={\Delta e_y}=0{,}001$, ${\Delta i}=0{,}002$
имеем $B_*<{(\Delta a)}^2$ и ${(\Delta e_x)}^2<{(\Delta a)}^2$. Формулы дают
\[
 J=0{,}002866974792,\qquad
 \varphi_+=165{,}02574460^\circ,\qquad
 \varphi_-=9{,}01738577^\circ,
\]
\begin {equation}
 \Delta\widehat{\boldsymbol V}_+=\begin {pmatrix}
 0{,}000084628214\\0{,}000755128527\\-0{,}000770943752
 \end {pmatrix},\qquad
 \Delta\widehat{\boldsymbol V}_-=\begin {pmatrix}
 -0{,}000139514328\\0{,}001244871473\\0{,}001270943752
 \end {pmatrix}.
 \label {eq:NDC-example}
\end {equation}
Выполнение сначала импульса с индексом $-$ даёт перелёт по меньшей
положительной дуге между указанными углами. При тех же $r_0$ и $\mu$
размерная стоимость составляет приблизительно $21{,}6343$ м/с.
Обратный порядок также допустим в идеальной задаче со свободным временем,
но задаёт другую пассивную дугу и другие ошибки при реализации
в точной модели.

\subsubsection{Вырожденный годограф-отрезок: ${\Delta e_y}=0$}
\label{sec:line}

Пусть ${\Delta i}>0$, ${\Delta e_y}=0$, ${(\Delta a)}^2\geqslant {(\Delta e_x)}^2$.
Оптимальный базис-вектор имеет вид
\begin {equation}
 J=\sqrt{{(\Delta i)}^2+{(\Delta a)}^2/4},\qquad
 \boldsymbol p(\varphi)=\frac{1}{J}
 \begin {pmatrix}0\\{\Delta a}/2\\{\Delta i}\cos\varphi\end {pmatrix}.
 \label {eq:line-primer}
\end {equation}
Годограф является отрезком, а не невырожденным эллипсом.
Поэтому вывод, использовавший деление на $R$, к нему неприменим.
Очевидно, $p^2=({(\Delta a)}^2/4+{(\Delta i)}^2\cos^2\varphi)/J^2\leqslant1$.
Два места приложения --- $0$ и $\pi$; при ${\Delta a}\ne0$
\begin {equation}
 \begin{gathered}
 \Delta\widehat V_1=\frac{J}{2}\left(1+\frac{\Delta e_x}{\Delta a}\right),\qquad
 \Delta\widehat V_2=\frac{J}{2}\left(1-\frac{\Delta e_x}{\Delta a}\right),\\[4pt]
 \Delta\widehat{\boldsymbol V}_1=\frac{\Delta\widehat V_1}{J}
 \begin{pmatrix}0\\\Delta a/2\\\Delta i\end{pmatrix},\qquad
 \Delta\widehat{\boldsymbol V}_2=\frac{\Delta\widehat V_2}{J}
 \begin{pmatrix}0\\\Delta a/2\\-\Delta i\end{pmatrix}.
 \end{gathered}
 \label {eq:line-impulses}
\end {equation}
Граничные условия проверяются прямой подстановкой, а неотрицательность
модулей эквивалентна $|{\Delta e_x}|\leqslant|{\Delta a}|$.
При ${\Delta a}={\Delta e_x}=0$ достаточно одного нормального импульса величины ${\Delta i}$ в узле.
Если ${\Delta e_y}=0$, но ${(\Delta a)}^2<{(\Delta e_x)}^2$, применяется NC, а не
\eqref{eq:line-impulses}.

В классификации Баранова годограф-отрезок выделяется как дополнительный
вырожденный тип. Его стоимость согласуется с пределом
\eqref{eq:NDC-cost}. Само наличие отдельной формулы не означает
разрыва оптимальной стоимости при ${\Delta e_y}\to0$.

% ==================== singular ====================
\subsubsection{Особый случай: постоянный модуль базис-вектора}
\label{sec:singular}

В особом случае модуль базис-вектора равен единице при любом
угле. Поэтому условия оптимальности определяют направления
импульсов, но не выделяют отдельные точки их приложения.
Для расчёта воспользуемся конструкцией Эдельбаума
\cite[с.~56]{edelbaum}: у всех импульсов выбирается одинаковое
значение $\sin^2(\varphi-\alpha)$, а величины определяются из оставшихся
граничных условий. Ниже эта конструкция доводится до явных формул
с неотрицательными модулями. Она использует не более трёх импульсов.
Достаточность двух импульсов, отдельно установленная
в~\cite[с.~55--56]{edelbaum}, не означает, что каждый удобный
расчётный алгоритм обязан использовать именно два.
Сначала считаем, что ${\Delta i}>0$. Тогда в рассматриваемом случае
${\Delta e_y}\ne0$. Введём
\begin {equation}
 \sigma=\operatorname {sgn}{\Delta e_y},\qquad
 J_s=\frac{1}{2}\sqrt{{(\Delta e_x)}^2+({|\Delta e_y|}+\sqrt3 {\Delta i})^2},\qquad
 \alpha=\operatorname {atan2}({\Delta e_y}+\sigma\sqrt3 {\Delta i},{\Delta e_x}).
 \label {eq:sc-price}
\end {equation}
Положим $\psi=\varphi-\alpha$ и рассмотрим базис-вектор
\begin {equation}
 \boldsymbol p(\varphi)=
 \begin {pmatrix}
  \tfrac{1}{2}\sin\psi\\
  \cos\psi\\
  -\sigma\tfrac{{\sqrt3}}{{2}}\sin\psi
 \end {pmatrix}.
 \label {eq:sc-primer}
\end {equation}
Он имеет требуемый вид~\eqref{eq:primer-components} и удовлетворяет
равенству $\lVert\boldsymbol p\rVert^2=
\tfrac{1}{4}\sin^2\psi+\cos^2\psi+\tfrac{3}{4}\sin^2\psi=1$.
Ему соответствуют множители
\begin {equation}
 \boldsymbol l=
 \begin {pmatrix}
  0\\ \tfrac{1}{2}\cos\alpha\\ \tfrac{1}{2}\sin\alpha\\
  -\sigma\tfrac{{\sqrt3}}{{2}}\cos\alpha\\
  \sigma\tfrac{{\sqrt3}}{{2}}\sin\alpha
 \end {pmatrix},
 \qquad
 \boldsymbol l^{\mathsf T}\boldsymbol b
 =\frac{1}{2}\bigl[{\Delta e_x}\cos\alpha+
       ({\Delta e_y}+\sigma\sqrt3 {\Delta i})\sin\alpha\bigr]=J_s.
 \label {eq:sc-certificate}
\end {equation}
Следовательно, $J_s$ является нижней оценкой стоимости любого
допустимого перехода. Чтобы доказать достижимость оценки, нужно
построить неотрицательные величины ${\Delta\widehat V_j}$ и углы $\varphi_j$, для которых
$\Delta\widehat{\boldsymbol V}_j={\Delta\widehat V_j}\boldsymbol p(\varphi_j)$ удовлетворяют всем пяти
граничным уравнениям. Именно эту задачу решаем далее.

\paragraph{Сведение граничных уравнений.}
Повернём компоненты изменения эксцентриситета на угол $\alpha$:
\begin {equation}
 {\Delta e'_x}={\Delta e_x}\cos\alpha+{\Delta e_y}\sin\alpha,
 \qquad {\Delta e'_y}=-{\Delta e_x}\sin\alpha+{\Delta e_y}\cos\alpha.
 \label {eq:sc-rotated-e}
\end {equation}
Штрих у $\Delta e'_x,\Delta e'_y$ обозначает компоненты в повёрнутых
осях, а не дифференцирование.
Подстановка~\eqref{eq:sc-primer} в уравнения перехода даёт
\begin {align}
 \sum_j {\Delta\widehat V_j}&=J_s,
 &\sum_j {\Delta\widehat V_j}\cos\psi_j&=\frac{{\Delta a}}{2},
 \label {eq:sc-moment-a}\\
 \sum_j {\Delta\widehat V_j}\sin^2\psi_j&=U,
 &\sum_j {\Delta\widehat V_j}\sin\psi_j\cos\psi_j&=V,
 \label {eq:sc-moment-b}
\end {align}
где
\begin {equation}
 U=\frac{2}{3}(2J_s-{\Delta e'_x})
   =\frac{{{\Delta i}({|\Delta e_y|}+\sqrt3 {\Delta i})}}{{\sqrt3 J_s}},
 \qquad
 V=\frac{2}{3}{\Delta e'_y}=-\frac{{\sigma {\Delta i} {\Delta e_x}}}{{\sqrt3 J_s}}.
 \label {eq:sc-uv}
\end {equation}
Например, первое уравнение для эксцентриситета в повёрнутых осях
имеет вид
\[
 {\Delta e'_x}=\sum_j{\Delta\widehat V_j}(\tfrac{1}{2}\sin^2\psi_j+2\cos^2\psi_j)
 =2J_s-\tfrac{3}{2}U.
\]
Второе уравнение даёт ${\Delta e'_y}=\tfrac{3}{2}V$.
Нормальные граничные условия в тех же осях записываются как
\[
 -\sigma\frac{{\sqrt3}}{{2}}U=-{\Delta i}\sin\alpha,
 \qquad
 -\sigma\frac{{\sqrt3}}{{2}}V={\Delta i}\cos\alpha.
\]
Они выполнены в силу~\eqref{eq:sc-uv}. Таким образом, четыре
соотношения~\eqref{eq:sc-moment-a}--\eqref{eq:sc-moment-b}
обеспечивают все пять требуемых изменений орбитальных параметров
и заданную суммарную стоимость.

\paragraph{Одинаковый квадрат синуса для всех импульсов.}
Положим $w_j={\Delta\widehat V_j}/J_s$; тогда $w_j\ge0$ и $\sum_jw_j=1$.
Следуя указанному построению Эдельбаума, потребуем
\begin {equation}
 \sin^2\psi_j=s_*^2,\qquad
 s_*^2=\frac{{U}}{{J_s}},\qquad c_*^2=1-s_*^2,
 \qquad s_*\ge0,\quad c_*\ge0.
 \label {eq:sc-equal-sine}
\end {equation}
Тогда условие на $U$ выполнено автоматически при любой разбивке
суммарного импульса. Для невырожденного случая $0<s_*<1$
введём угол $\beta=\operatorname {atan2}(s_*,c_*)\in(0,\pi/2)$.
Имеются четыре возможных положения; соответствующие знаки
косинуса и произведения синуса на косинус приведены в таблице:
\begin{center}
\begin{tabular}{c|c|cc}
 $j$&$\psi_j$&$\cos\psi_j/c_*$&
 $\sin\psi_j\cos\psi_j/(s_*c_*)$\\\hline
 1&$\pi-\beta$&$-1$&$-1$\\
 2&$-\beta$&$1$&$-1$\\
 3&$\beta$&$1$&$1$\\
 4&$\pi+\beta$&$-1$&$1$
\end{tabular}
\end{center}
Два ещё не выполненных условия записываются в виде
\begin {equation}
 -w_1+w_2+w_3-w_4=\xi,\qquad
 -w_1-w_2+w_3+w_4=\zeta,
 \label {eq:sc-weight-balances}
\end {equation}
где
\begin {equation}
 \xi=\frac{{{\Delta a}}}{{2J_sc_*}},\qquad
 \zeta=\frac{{V}}{{J_ss_*c_*}}.
 \label {eq:sc-rectangle-point}
\end {equation}
Каждому из четырёх углов отвечает одна вершина квадрата
с координатами $(\pm1,\pm1)$.
Следовательно, найти неотрицательные веса означает представить
точку $(\xi,\zeta)$ как взвешенное среднее этих вершин.
Это возможно тогда и только тогда, когда
$|\xi|\le1$ и $|\zeta|\le1$.

\begin{figure}[htbp]
\centering
\begin{tikzpicture}[scale=1.6,>=Latex]
 \fill[accent!8] (-1,-1)--(1,-1)--(1,1)--cycle;
 \draw[->,gray] (-1.35,0)--(1.42,0) node[right] {$\xi$};
 \draw[->,gray] (0,-1.35)--(0,1.42) node[above] {$\zeta$};
 \draw[thick] (-1,-1) rectangle (1,1);
 \draw[accent] (-1,-1)--(1,1);
 \foreach \xx/\yy in {-1/-1,1/-1,1/1,-1/1}
  {\fill (\xx,\yy) circle(.9pt);}
 \node[below left] at (-1,-1) {$1$};
 \node[below right] at (1,-1) {$2$};
 \node[above right] at (1,1) {$3$};
 \node[above left] at (-1,1) {$4$};
 \fill[accent] (.74,-.12) circle(1.2pt);
 \node[below,accent] at (.63,-.2) {$(\xi,\zeta)$};
\end{tikzpicture}
\caption{Выбор трёх импульсов в конструкции Эдельбаума.
При $\zeta\leqslant\xi$ используются вершины $1,2,3$;
при $\zeta\geqslant\xi$ --- вершины $1,3,4$.
Веса являются долями суммарной характеристической скорости.}
\label{fig:SC-square}
\end{figure}

Покажем, что эти ограничения совпадают с областью особого решения.
Из~\eqref{eq:sc-uv} следует
\begin {align}
 c_*^2&=\frac{{{(\Delta e_x)}^2+{(\Delta e_y)}^2+2{\Delta i}{|\Delta e_y|}/\sqrt3-{(\Delta i)}^2}}{{4J_s^2}},
 \label {eq:sc-cos-square}\\
 s_*^2c_*^2-\left(\frac{{V}}{{J_s}}\right)^2
 &=\frac{{{\Delta i}(\sqrt3{|\Delta e_y|}-{\Delta i})}}{{3J_s^2}}.
 \label {eq:sc-rectangle-determinant}
\end {align}
Поэтому $|\xi|\le1$ и $|\zeta|\le1$ эквивалентны
\begin {equation}
 \boxed{\qquad
 {(\Delta i)}^2\le3{(\Delta e_y)}^2,\qquad
 {(\Delta a)}^2\le {(\Delta e_x)}^2+{(\Delta e_y)}^2+\frac{{2{\Delta i}|{\Delta e_y}|}}{{\sqrt3}}-{(\Delta i)}^2.
 \qquad}
 \label {eq:sc-region}
\end {equation}
Внутри этой области $c_*^2\ge0$, а $s_*^2\ge0$ следует
из~\eqref{eq:sc-uv}; при этом их сумма равна единице.
Таким образом, выбранные углы существуют.

\paragraph{Выбор трёх углов и величин импульсов.}
Диагональ $\zeta=\xi$ делит квадрат на два треугольника.
Если $\zeta\le\xi$, выбираем вершины $1,2,3$ и полагаем $w_4=0$.
Решение двух уравнений~\eqref{eq:sc-weight-balances} и условия
$w_1+w_2+w_3=1$ даёт
\begin {equation}
 w_1=\frac{{1-\xi}}{{2}},\qquad
 w_2=\frac{{\xi-\zeta}}{{2}},\qquad
 w_3=\frac{{1+\zeta}}{{2}},\qquad w_4=0.
 \label {eq:sc-weights-lower}
\end {equation}
Если $\zeta\ge\xi$, выбираем вершины $1,3,4$:
\begin {equation}
 w_1=\frac{{1-\zeta}}{{2}},\qquad w_2=0,\qquad
 w_3=\frac{{1+\xi}}{{2}},\qquad
 w_4=\frac{{\zeta-\xi}}{{2}}.
 \label {eq:sc-weights-upper}
\end {equation}
Все веса неотрицательны, поскольку $\xi,\zeta\in[-1,1]$
и выбор треугольника согласован со знаком $\xi-\zeta$.
Их сумма равна единице. На диагонали оба варианта совпадают
и ненулевых весов не более двух.
Можно пользоваться одной объединённой записью:
\begin {align}
 w_1&=\tfrac{1}{2}\bigl(1-\max(\xi,\zeta)\bigr),&
 w_2&=\tfrac{1}{2}\max(\xi-\zeta,0),\notag\\
 w_3&=\tfrac{1}{2}\bigl(1+\min(\xi,\zeta)\bigr),&
 w_4&=\tfrac{1}{2}\max(\zeta-\xi,0).
 \label {eq:sc-weights-unified}
\end {align}
Один из весов $w_2,w_4$ всегда равен нулю; поэтому алгоритм
действительно требует не более трёх импульсов.

Окончательно
\begin {equation}
 \varphi_j=\alpha+\psi_j\pmod{2\pi},\qquad
 {\Delta\widehat V_j}=J_sw_j,\qquad
 \Delta\widehat{\boldsymbol V}_j={\Delta\widehat V_j}\boldsymbol p(\varphi_j).
 \label {eq:sc-impulses}
\end {equation}
Члены с $w_j=0$ исключаются. Углы следует расположить в требуемом
порядке движения; номера вершин квадрата сами по себе не являются
порядковыми номерами манёвров по времени.
Все четыре соотношения~\eqref{eq:sc-moment-a}--\eqref{eq:sc-moment-b}
выполнены по построению, а значит, выполнены и пять граничных
уравнений исходной задачи. Кроме того,
$\sum_j{\Delta\widehat V_j}=J_s$ и $\lVert\boldsymbol p\rVert=1$.
Нижняя оценка~\eqref{eq:sc-certificate} достигнута, что доказывает
глобальную оптимальность в линейной модели.

\paragraph{Граничные случаи и порядок расчёта.}
При $c_*=0$ формулы для $\xi$ и $\zeta$ неприменимы.
Однако из~\eqref{eq:sc-cos-square} и условий области следуют
${\Delta a}={\Delta e_x}=0$, ${\Delta i}=\sqrt3{|\Delta e_y|}$, а из~\eqref{eq:sc-uv} --- $V=0$.
Достаточен один импульс с ${\Delta\widehat V_1}=J_s$ при $\psi_1=\pi/2$.
Он удовлетворяет всем четырём приведённым граничным условиям.
При $s_*=0$ имеем ${\Delta i}=0$, и следует применять компланарное решение.
Если $J_s=0$, то в области SC все изменения равны нулю и импульсы
не требуются.

Практический расчёт выполняется в следующем порядке:
\begin{enumerate}
 \item Проверить область~\eqref{eq:sc-region}; вычислить $J_s$ и $\alpha$.
 \item Найти $U,V$ по~\eqref{eq:sc-uv} и $s_*,c_*$
       по~\eqref{eq:sc-equal-sine}; отдельно обработать нулевые значения.
 \item Найти $\beta$, четыре угла из таблицы и координаты
       $(\xi,\zeta)$ по~\eqref{eq:sc-rectangle-point}.
 \item Вычислить веса по~\eqref{eq:sc-weights-unified}, исключить
       нулевые импульсы и восстановить векторы~\eqref{eq:sc-impulses}.
 \item Проверить пять исходных граничных уравнений и сумму модулей.
\end{enumerate}
Вблизи границы $c_*=0$ для вычисления его квадрата удобно использовать
алгебраически равносильную запись
\[
 c_*^2=\frac{{{(\Delta e_x)}^2+(\sqrt3{|\Delta e_y|}-{\Delta i})({|\Delta e_y|}/\sqrt3+{\Delta i})}}{{4J_s^2}}.
\]
Она не требует вычитания близких положительных слагаемых.
Небольшой выход вычисленного $\xi$ или $\zeta$ за $[-1,1]$
можно исправлять только в пределах ошибки округления после
проверки принадлежности области; значительный выход означает ошибку.

\paragraph{Численный пример.}
Пусть
\[
 {\Delta a}=10^{-3},\qquad {\Delta e_x}=0,\qquad {\Delta e_y}=10^{-3},\qquad {\Delta i}=2\cdot10^{-4}.
\]
Оба неравенства~\eqref{eq:sc-region} выполнены. Получаем
$J_s=6{,}73205081\cdot10^{-4}$, $\alpha=\pi/2$,
$\beta=35{,}85253018^\circ$ и $\zeta=0$.
Результат после сортировки углов приведён в таблице.
Компоненты и модули указаны в единицах $10^{-4}$:
\begin{center}
\begin{tabular}{c|rrrr}
 $\varphi_j$, град & ${\Delta\widehat V_{rj}}$ & ${\Delta\widehat V_{tj}}$ & ${\Delta\widehat V_{zj}}$ & ${\Delta\widehat V_j}$\\\hline
 $54{,}14746982$ & $-0{,}90327174$ & $2{,}50000000$ & $1{,}56451254$ & $3{,}08441232$\\
 $125{,}85253018$ & $0{,}98574228$ & $2{,}72825506$ & $-1{,}70735571$ & $3{,}36602540$\\
 $234{,}14746982$ & $0{,}08247054$ & $-0{,}22825506$ & $-0{,}14284316$ & $0{,}28161309$
\end{tabular}
\end{center}
Подстановка неокруглённых значений даёт
\[
 \sum_jB(\varphi_j)\Delta\widehat{\boldsymbol V}_j
 =\begin {pmatrix}10^{-3}&0&10^{-3}&0&2\cdot10^{-4}\end {pmatrix}^{\mathsf T},
 \qquad \sum_j{\Delta\widehat V_j}=6{,}73205081\cdot10^{-4}.
\]
Для перевода импульсов в размерные единицы все компоненты
таблицы необходимо умножить на $10^{-4}V_0$.

% ==================== validation ====================
\subsubsection{Сводный порядок расчёта пространственного перехода}

Для реализации не требуется решать общую нелинейную задачу оптимизации
по всем углам и компонентам одновременно. Достаточно следующего порядка.
\begin{enumerate}
 \item Построить общие оси и вычислить ${\Delta a},{\Delta e_x},{\Delta e_y},{\Delta i}$.
 При ${\Delta i}=0$ перейти к компланарному алгоритму.
 \item При ${(\Delta a)}^2\leqslant {(\Delta e_x)}^2$ и ${(\Delta i)}^2\geqslant3{(\Delta e_y)}^2$ использовать NC.
 Обработать ${\Delta e_x}=0$ отдельно, как указано в разделе~\ref{sec:NC}.
 \item Иначе, при ${(\Delta i)}^2\leqslant3{(\Delta e_y)}^2$ и ${(\Delta a)}^2\leqslant B_*$
 использовать SC с явным выбором не более трёх ненулевых весов.
 \item Для оставшихся данных использовать NDC.
 При ${\Delta e_y}=0$ применить формулы годографа-отрезка
 \eqref{eq:line-impulses}; при ${\Delta e_y}\ne0$ ---
 \eqref{eq:NDC-eta}--\eqref{eq:NDC-magnitudes}.
 \item Исключить нулевые импульсы, привести углы к одному обороту
 и выбрать их хронологический порядок.
 Восстановить размерные компоненты умножением на $V_0$.
\end{enumerate}
Порядок проверок задаёт однозначную вычислительную ветвь на общих
границах, хотя само оптимальное управление там может быть неединственным.
При численном сравнении учитывают масштаб данных и точность вычислений.
Отрицательный модуль импульса нельзя исправлять изменением его знака:
это изменяет направление относительно базис-вектора и, как правило,
нарушает граничные условия или оптимальность.

Для всех полученных решений обязательна проверка в исходных,
а не в промежуточно преобразованных уравнениях:
\begin {equation}
 \boldsymbol r_b=\sum_jB(\varphi_j)\Delta\widehat{\boldsymbol V}_j-
 \begin {pmatrix}{\Delta a}\\{\Delta e_x}\\{\Delta e_y}\\0\\{\Delta i}\end {pmatrix},\qquad
 r_J=\sum_j\lVert\Delta\widehat{\boldsymbol V}_j\rVert-J_*.
 \label {eq:algorithm-residuals}
\end {equation}
В точной арифметике обе невязки равны нулю.
Дополнительно проверяются неотрицательность модулей, направления
импульсов по базис-вектору и его глобальное ограничение по норме.
Для NC ограничение следует из \eqref{eq:NC-primer-norm},
для NDC --- из вогнутого квадратичного выражения
\eqref{eq:NDC-normalization}, для SC --- из тождества $p^2=1$.
Проверка нормы только на конечной сетке углов сама по себе не заменяет
эти аналитические условия.

\subsection{Проверка в точной кеплеровой модели}
\label{sec:exact-verification}

\subsubsection{Что именно проверяется}

Удовлетворение пяти линейным условиям означает точность первого
порядка по малым параметрам. Для оценки фактического выхода на орбиту
нужно выполнить импульсы в нелинейной модели и заново вычислить
конечные элементы. После каждого импульса положение сохраняется,
а скорость заменяется на $\boldsymbol V^+=\boldsymbol V^-+
\Delta\boldsymbol V$. Между импульсами используются уравнения
\eqref{eq:cartesian-dynamics} при нулевом управлении.

По состоянию после последнего импульса вычисляют
\begin {align}
 a_{\rm calc}&=\left(\frac{{2}}{{\lVert\boldsymbol R\rVert}}
                   -\frac{{\lVert\boldsymbol V\rVert^2}}{{\mu}}\right)^{-1},
 \label {eq:exact-a-check}\\
 \boldsymbol e_{\rm calc}
 &=\frac{{\boldsymbol V\times(\boldsymbol R\times\boldsymbol V)}}{{\mu}}
          -\frac{{\boldsymbol R}}{{\lVert\boldsymbol R\rVert}},\qquad
 \boldsymbol k_{\rm calc}
 =\frac{{\boldsymbol R\times\boldsymbol V}}{{\lVert\boldsymbol R\times\boldsymbol V\rVert}}.
 \label {eq:exact-ei-check}
\end {align}
Характеристики ошибки удобно выбирать без деления на малые
эксцентриситеты и наклонения:
\begin {equation}
 \begin {aligned}
 \delta a&=a_{\rm calc}-a_f,\\
 \delta\boldsymbol e_\parallel
 &=\begin {pmatrix}
 (\boldsymbol e_{\rm calc}-\boldsymbol e_f)\cdot\boldsymbol e_X\\
 (\boldsymbol e_{\rm calc}-\boldsymbol e_f)\cdot\boldsymbol e_Y
 \end {pmatrix},\\
 \delta\gamma&=\operatorname {atan2}
 \bigl(\lVert\boldsymbol k_{\rm calc}\times\boldsymbol k_f\rVert,
       \boldsymbol k_{\rm calc}\cdot\boldsymbol k_f\bigr).
 \end {aligned}
 \label {eq:exact-residual-measures}
\end {equation}
Погрешность аргумента перицентра для почти круговой орбиты может быть
большой даже при малой ошибке вектора эксцентриситета. Поэтому она
не является надёжным самостоятельным показателем точности такого перехода.

\subsubsection{Воспроизводимая схема численного опыта}

Для примеров предыдущих подразделов примем одинаковые
$r_0=7000$ км и $\mu=398600{,}4418$ км$^3$/с$^2$.
Начальная орбита круговая, лежит в плоскости $XY$ и имеет радиус
$a_i=r_0(1-{\Delta a}/2)$; конечная большая полуось равна $a_f=r_0(1+{\Delta a}/2)$.
Конечная плоскость получена поворотом начальной на угол ${\Delta i}$ вокруг
оси $X$. Для неё зададим точный вектор эксцентриситета
\[
 \boldsymbol e_f={\Delta e_x}\boldsymbol e_X+{\Delta e_y}\boldsymbol e_Y
                      +{\Delta e_y}\tan {\Delta i}\,\boldsymbol e_Z.
\]
Он лежит в конечной плоскости и имеет заданные проекции $({\Delta e_x},{\Delta e_y})$.
Такое задание определяет точные орбиты, с которыми сравнивается результат.

Алгоритмические углы приводятся к $[0,2\pi)$ и сортируются по
возрастанию. Начальный аппарат устанавливается на своей круговой
орбите в первом из этих углов. Между последовательными импульсами
задаётся время
\begin {equation}
 \Delta t_j=\frac{{\varphi_{j+1}-\varphi_j}}{{n_0}}.
 \label {eq:verification-timing}
\end {equation}
Это время по опорному кругу; в точном движении фактическая долгота
к следующему манёвру немного отличается от алгоритмической.
Векторы импульсов прикладываются с рассчитанными компонентами
в цилиндрическом базисе \eqref{eq:cylindrical-basis}, построенном
по фактическому положению аппарата. Таким образом, способ переноса
линейного решения в точную модель определён полностью.
Другой способ, например включение при достижении заданной геометрической
долготы, даст другие члены второго порядка и должен проверяться отдельно.

Пассивные дуги интегрировались в единицах $r_0,V_0,1/n_0$
методом DOP853 с относительным допуском $2\cdot10^{-13}$ и абсолютным
допуском $2\cdot10^{-15}$. При этих допусках численная ошибка существенно
меньше приведённых ниже невязок модели. Результаты относятся к
неокруглённым импульсам из аналитических формул.

\begin{table}[htbp]
\centering
\caption{Невязки конечной орбиты при выполнении линейного решения в точной модели.
$\delta\gamma$ выражена в угловых секундах.}
\label{tab:exact-check}
\renewcommand{\arraystretch}{1.2}
\begin{tabular}{@{}lrrrr@{}}
\toprule
Случай & $\Delta V_\Sigma$, м/с & $\delta a$, м &
$\lVert\delta\boldsymbol e_\parallel\rVert$ & $\delta\gamma$, угл.~с\\
\midrule
Компланарный & $8{,}43674$ & $17{,}5107$ & $3{,}28783\cdot10^{-6}$ & $0$\\
NC & $17{,}29018$ & $43{,}7774$ & $1{,}95361\cdot10^{-6}$ & $0{,}33260$\\
NDC & $21{,}63434$ & $35{,}7720$ & $2{,}77055\cdot10^{-6}$ & $0{,}46261$\\
SC & $5{,}08004$ & $4{,}1859$ & $5{,}04182\cdot10^{-7}$ & $0{,}02701$\\
\bottomrule
\end{tabular}
\end{table}

Линейные граничные невязки для этих примеров менее $10^{-17}$
в безразмерной норме. Ненулевые значения в таблице~\ref{tab:exact-check}
возникают при переходе к точной модели.
Они не противоречат доказанной линейной оптимальности.
Значение $\Delta V_\Sigma$ при указанной реализации сохраняется:
компоненты тех же импульсов пересчитываются между ортонормированными
базисами. Однако полученная точная орбита отличается от заданной,
и расход после устранения этой невязки заранее неизвестен.

Проверим порядок ошибки, уменьшая одновременно ${\Delta a},{\Delta e_x},{\Delta e_y},{\Delta i}$ примера NDC
в $1/\kappa$ раз. Углы импульсов при таком масштабировании не меняются,
а их величины пропорциональны $\kappa$.
Введём безразмерную норму пяти конечных ошибок
\[
 \varepsilon_b=\left\lVert
 \left(\delta a/r_0,\ \delta e_x,\ \delta e_y,
             \delta i_y,\ \delta i_x\right)^{\mathsf T}\right\rVert.
\]
Компоненты наклонения для этой проверки вычисляются из точной
нормали как $i_x=i\cos\Omega$, $i_y=i\sin\Omega$ в опорном базисе.
Получаем
\begin{center}
\begin{tabular}{@{}crr@{}}
\toprule
$\kappa$ & $\varepsilon_b$ & Отношение к предыдущей невязке\\
\midrule
$1$ & $6{,}23065\cdot10^{-6}$ & ---\\
$1/2$ & $1{,}55762\cdot10^{-6}$ & $0{,}2500$\\
$1/4$ & $3{,}89401\cdot10^{-7}$ & $0{,}2500$\\
$1/8$ & $9{,}73497\cdot10^{-8}$ & $0{,}2500$\\
\bottomrule
\end{tabular}
\end{center}
Уменьшение входных отклонений вдвое уменьшает невязку приблизительно
вчетверо, что согласуется со вторым порядком отброшенных членов.
Это проверка конкретной серии примеров при ограниченной длительности,
а не универсальная оценка ошибки для любых близких орбит и любых
добавленных оборотов.

\subsubsection{Уточнение и предельные проверки}

Если фактическая невязка неприемлема, линейное решение служит начальным
приближением для коррекции параметров манёвров в точной модели.
Пусть $\boldsymbol z$ содержит свободные компоненты импульсов
и выбранные параметры времени, а $\boldsymbol F(\boldsymbol z)$ ---
пять точных невязок орбиты. Первый шаг дифференциальной коррекции
определяется системой
\[
 \frac{{\partial\boldsymbol F}}{{\partial\boldsymbol z}}
 \delta\boldsymbol z=-\boldsymbol F(\boldsymbol z).
\]
При избытке переменных необходимо дополнительно задать правило выбора
поправки или решать задачу минимизации стоимости с точными ограничениями.
Само обнуление невязки доказывает достижение конечной орбиты,
но не доказывает оптимальность уточнённого перехода.
Задача такого уточнения выходит за пределы аналитического первого порядка.

Несколько простых пределов полезны для проверки нормировки всех формул:
\begin {align*}
 {\Delta i}={\Delta e_x}={\Delta e_y}=0&:\qquad J_*=|{\Delta a}|/2,\\
 {\Delta a}={\Delta i}=0&:\qquad J_*=\sqrt{{(\Delta e_x)}^2+{(\Delta e_y)}^2}/2,\\
 {\Delta a}={\Delta e_x}={\Delta e_y}=0&:\qquad J_*={\Delta i},\\
 {\Delta e_x}={\Delta e_y}=0&:\qquad J_*=\sqrt{{(\Delta i)}^2+{(\Delta a)}^2/4}.
\end {align*}
Кроме того, оптимальная стоимость должна быть положительно однородной:
$J_*(\kappa\boldsymbol b)=\kappa J_*(\boldsymbol b)$ при $\kappa\geqslant0$.
Эти свойства непосредственно следуют из линейных граничных условий
и используются наряду с подстановкой полного решения.
Именно они позволяют обнаруживать ошибочные коэффициенты при
переносе формул между различными обозначениями источников.

В алгоритмах раздела~\ref{sec:spatial} радиальная компонента не запрещена.
Описанные Барановым решения с дополнительным условием отсутствия
радиальных импульсов и его универсальная приближённая схема решают
отдельную ограниченную или приближённую задачу.
Их нельзя подставлять вместо общего оптимального решения без указания
этого дополнительного предположения.

% ==================== back ====================
\clearpage
\begin{thebibliography}{9}
\addcontentsline{toc}{section}{Литература}
\bibitem{lawden}
Lawden D.~F. \textit{Optimal Trajectories for Space Navigation}.
London: Butterworths, 1963. Гл.~3, \S~5.1, гл.~6.
\bibitem{lawden-ru}
Лоуден Д.~Ф. \textit{Оптимальные траектории для космической навигации}.
Пер. с англ. В.~К.~Исаева. М.: Мир, 1966.
\bibitem{edelbaum}
Edelbaum T.~N. A General Solution for Minimum Impulse Transfers
in the Near Vicinity of a Circular Orbit //
\textit{First Compilation of Papers on Trajectory Analysis and Guidance Theory}.
NASA SP-141, 1967. P.~37--71.
\href{https://ntrs.nasa.gov/citations/19670020044}{Запись NASA NTRS}.
\bibitem{gobetz}
Gobetz F.~W., Washington M., Edelbaum T.~N.
\textit{Minimum-Impulse Time-Free Transfer Between Elliptic Orbits}.
NASA CR-636, 1966. Линейная задача: с.~25--42;
приложение~C: с.~45--49.
\href{https://ntrs.nasa.gov/citations/19660030604}{Запись NASA NTRS}.
\bibitem{baranov-ru}
Баранов А.~А. \textit{Разработка методов расчета параметров маневров
космических аппаратов в окрестностях круговой орбиты}.
Диссертация на соискание учёной степени доктора физико-математических наук.
М.: ИПМ им.~М.~В.~Келдыша РАН, 2018.
Электронная публикация 2019 г. Гл.~II, с.~50--79.
\bibitem{baranov-en}
Baranov A. \textit{Spacecraft Manoeuvring in the Vicinity of a Near-Circular Orbit}.
Newcastle upon Tyne: Cambridge Scholars Publishing, 2022.
Chapter~2, p.~41--74.
\end{thebibliography}
\end{document}
